% -----------------------------------------------------------
\section{Epistemology of revelation}
% -----------------------------------------------------------
	\label{EPIST-REVELATION}
	
	This entry is for the theory and epistemology of revelation. For uses of the word ``revelation,'' see \hyperref[REVELATION]{Revelation}.
	
% ----------------------
\subsection{See also}
% ----------------------

	\hyperref[GOODNESS-truth]{Goodness and truth},

% % ----------------------
% \subsection{1828 American Dictionary of the English Language} % *1828
% % ----------------------
	% \label{EPIST-REVELATION-1828}

	% \begin{compactitem}
		% \item
	% \end{compactitem}

% % ----------------------
% \subsection{Oxford English Dictionary} % *OED
% % ----------------------
	% \label{EPIST-REVELATION-OED}

	% \begin{compactitem}
		% \item[]
	% \end{compactitem}

% ----------------------
\subsection{Scriptural Usage} % *SCRIPTURES
% ----------------------
	\label{EPIST-REVELATION-SCRIPTURES}

	% -----
	\subsubsection{Old Testament} % *OT
	% -----
		\label{EPIST-REVELATION-OT}
	
		% --
		\paragraph{KJV}
		% --
			\label{EPIST-REVELATION-OT-KJV}
	
			% \begin{tabular}{ll}
				% Total occurrences in the OT & 
			% \end{tabular}
			
			\begin{compactdesc}
				\item[{\hyperref[DEUTERONOMY-8-3]{Deuteronomy 8:3}}]
			\end{compactdesc}
			
		% % --
		% \paragraph{Hebrew Bible}
		% % --
			% \label{EPIST-REVELATION-HEBREW}
		
			% %
			% \subparagraph{\texorpdfstring{\he{sample word}}{}}
			% %
				% \label{PAGA} % whatever is in step bible without periods
			
				% \begin{compactdesc}
					% \item[HALOT , PDF ] \textbf{}
						% \begin{compactitem}
							% \item[1]
						% \end{compactitem}
					% \item[BDB , PDF ] \textbf{}
						% \begin{compactitem}
							% \item[1]
						% \end{compactitem}
					% \item[DCH PDF ] \textbf{}
						% \begin{compactitem}
							% \item[1]
						% \end{compactitem}
				% \end{compactdesc}
				
				% \begin{compactdesc}
					% \item[Some OT uses] \textbf{}
						% \begin{compactdesc}
							% \item[]
						% \end{compactdesc}
				% \end{compactdesc}
			
		% % --
		% \paragraph{Septuagint}
		% % --
			% \label{EPIST-REVELATION-LXX}
		
			% %
			% \subparagraph{\texorpdfstring{\gr{sample word}}{}}
			% %
		
	% -----
	\subsubsection{New Testament} % *NT
	% -----
		\label{EPIST-REVELATION-NT}
	
		% --
		\paragraph{KJV}
		% --
			\label{EPIST-REVELATION-NT-KJV}
			
	
			% \begin{tabular}{ll}
				% Total occurrences in the NT & 
			% \end{tabular}
			
			\begin{compactdesc}
				\item[{\hyperref[MATTHEW-4-4]{Matthew 4:4}}]
				\item The ``good tree'' verses from the sermon on the mount
				\item[{\hyperref[MATTHEW-7-11]{Matthew 7:7--11}}]
				\item[{\hyperref[MATTHEW-7-15]{Matthew 7:15--20}}]
				\item[{\hyperref[MATTHEW-9-16]{Matthew 9:16--17}}] \phantom{.}
					\begin{compactitem}
						\item These are the ``new bottles'' verses. See also \hyperref[MARK-2-18]{Mark 2:18--22} and \hyperref[LUKE-5-33]{Luke 5:33--39}. Verse 39 in Luke is particularly important. 
					\end{compactitem}
				\item[{\hyperref[MATTHEW-10-19]{Matthew 10:19--20}}]
				\item[{\hyperref[MATTHEW-10-26]{Matthew 10:26--27}}]
				\item[{\hyperref[MATTHEW-11-19]{Matthew 11:19}}]
				\item[{\hyperref[MATTHEW-11-25]{Matthew 11:25--27}}]
				\item[{\hyperref[MATTHEW-12-33]{Matthew 12:33--37}}]
				\item[{\hyperref[MATTHEW-13-1]{Matthew 13:1--9}}] \phantom{.}
					\begin{compactitem}
						\item This is the parable of the sower---good for the principle of context. The ground matters. See also \hyperref[MARK-4-3]{Mark 4:3--9} and \hyperref[LUKE-8-4]{Luke 8:4--8}. Think also about Alma 32 in light of these verses.
					\end{compactitem}
				\item[{\hyperref[MATTHEW-13-10]{Matthew 13:10--17}}] \phantom{.}
					\begin{compactitem}
						\item This is what Christ says in response to the question, why do you speak to the people in parables? Notice the key verses 11--12 in particular.
					\end{compactitem}
				\item[{\hyperref[MATTHEW-13-34]{Matthew 13:34--35}}] 
				\item[{\hyperref[MATTHEW-16-13]{Matthew 16:13--20}}] 
				\item[{\hyperref[MATTHEW-23-34]{Matthew 23:34}}] \phantom{.}
					\begin{compactitem}
						\item Include philosophers on the list of people God sends to warn others? Maybe a bit of a stretch.
					\end{compactitem}
				\item[{\hyperref[MATTHEW-24-45]{Matthew 24:45}}] \phantom{.}
					\begin{compactitem}
						\item Reminds me of the distinction between milk and meat, and that makes me think of the verses in \hyperref[ALMA-12]{Alma 12} that talk about certain people receiving more and only being allowed to share a small amount.
					\end{compactitem}
				\item[{\hyperref[MATTHEW-25-14]{Matthew 25:14--30}}] \phantom{.}
					\begin{compactitem}
						\item With the different amounts of ``talents'' you could see something like the principle of \hyperref[ALMA-12]{Alma 12} here too. See also \hyperref[DC82-3]{DC 82:3}.
						\item And don't forget about verse 29, either, which is a broad principle that can apply to epistemological stuff.
					\end{compactitem}
				\item[{\hyperref[MARK-4-3]{Mark 4:3--34}}] \phantom{.}
					\begin{compactitem}
						\item There's a huge amount here---it's important to break this down into its various divisions and relate it to other verses, like \hyperref[ALMA-32]{Alma 32}.
						\item Verse 11 may be a reference to \hyperref[ALMA-12-9]{Alma 12:9--11}.
					\end{compactitem}
				\item[{\hyperref[MARK-6-51]{Mark 6:51--52}}] \phantom{.}
					\begin{compactitem}
						\item I'm specifically thinking here of how they were not considering one event, or weren't considering it enough, and this affected how they responded to another event.
					\end{compactitem}
				\item[{\hyperref[MARK-8-22]{Mark 8:22--26}}] \phantom{.}
					\begin{compactitem}
						\item A very interesting story to think about what Christ did to the man the first time he healed him, and what changes occurred the second time. Anything with concepts or something like that going on?
					\end{compactitem}
				\item[{\hyperref[MARK-13-11]{Mark 13:11}}]
				\item[{\hyperref[LUKE-2-25]{Luke 2:25--27}}] \phantom{.}
					\begin{compactitem}
						\item You could consider up to and incluidng verse 35 to be relevant here too.
					\end{compactitem}
				\item[{\hyperref[LUKE-8-4]{Luke 8:4--18}}] \phantom{.}
					\begin{compactitem}
						\item This is Luke's version of the parable of the sower.
						\item I wonder if \hyperref[LUKE-9-44]{Luke 9:44--45} is a clue about how to understand the bit here about ``hearing they don't hear.'' Not sure though. 
					\end{compactitem}
				\item[{\hyperref[LUKE-9-44]{Luke 9:44--45}}] \phantom{.}
				\item[{\hyperref[LUKE-10-21]{Luke 10:21--24}}] \phantom{.}
				\item[{\hyperref[LUKE-11-46]{Luke 11:46--52}}] \phantom{.}
					\begin{compactitem}
						\item You might want to expand this to the verses before 46 talking about the scribes and pharisees. In any case, the key verse is 52, but that has to be read in light of the ones before it and in light of the entire passage.
					\end{compactitem}
				\item[{\hyperref[LUKE-12-1]{Luke 12:1--3}}] \phantom{.}
				\item[{\hyperref[LUKE-12-11]{Luke 12:11--12}}] \phantom{.}
				\item[{\hyperref[LUKE-16-10]{Luke 16:10}}] \phantom{.}
					\begin{compactitem}
						\item This is part of an unusual parable, but the principle in this verse is particularly important.
					\end{compactitem}
				\item[{\hyperref[LUKE-18-34]{Luke 18:34}}] \phantom{.}
					\begin{compactitem}
						\item Requires a bit of context, but a very interesting example to look at.
					\end{compactitem}
				\item[{\hyperref[LUKE-19-12]{Luke 19:12--27}}] \phantom{.}
					\begin{compactitem}
						\item Not just about knowledge, also about other things.
					\end{compactitem}
				\item[{\hyperref[LUKE-21-14]{Luke 21:14--15}}] \phantom{.}
				\item[{\hyperref[LUKE-24-16]{Luke 24:16}}] \phantom{.}
					\begin{compactitem}
						\item This verse is part of the story of the two men who don't recognize Jesus after the resurrection. See verses 31--32 for the culmination.
					\end{compactitem}
				\item[{\hyperref[LUKE-24-44]{Luke 24:44--45}}] \phantom{.}
				\item[{\hyperref[LUKE-24-49]{Luke 24:49}}] \phantom{.}
				\item[{\hyperref[JOHN-3-34]{John 3:34}}] \phantom{.}
				\item[{\hyperref[JOHN-4-39]{John 4:39--42}}] \phantom{.}
				\item[{\hyperref[JOHN-5-17]{John 5:17--47}}] \phantom{.}
					\begin{compactitem}
						\item Not every verse in this block is relevant, but many of them are.
					\end{compactitem}
				\item[{\hyperref[JOHN-7-16]{John 7:16--18}}] \phantom{.}
				\item[{\hyperref[JOHN-15-26]{John 15:26}}] \phantom{.}
				\item[{\hyperref[JOHN-16]{John 16}}] \phantom{.}
					\begin{compactitem}
						\item Not necessarily all of the verses in the chapter, but many of them have to do with epistemology, including the relationship between Christ's presence and the Comforter, and so on.
						\item See 13--16, for example.
					\end{compactitem}
				\item[{\hyperref[JOHN-20-5]{John 20:5--18}}] \phantom{.}
					\begin{compactitem}
						\item A few points here. First, Mary looks into the sepulchre and sees two angels that John and Peter apparently don't see.
						\item The second is that Mary doesn't recognize the Lord when she first sees him. Why not? And note that the answer might be as simple as, she's distracted and upset, and not looking at him carefully when he's talking (she ``saw Jesus standing'' in verse 14 though).
					\end{compactitem}
				\item[{\hyperref[JOHN-21-4]{John 21:4--12}}] \phantom{.}
					\begin{compactitem}
						\item This story is a bit similar to Mary's in the previous chapter.
					\end{compactitem}
				\item[{\hyperref[ROMANS-3-5]{Romans 3:5}}] \phantom{.}
					\begin{compactitem}
						\item I'm specifically thinking of the ``I speak as a man'' part, though I'm not sure that it's necessarily a constructive/contextual thing.
					\end{compactitem}
				\item[{\hyperref[ROMANS-3-20]{Romans 3:20}}]
				\item[{\hyperref[ROMANS-6-19]{Romans 6:19}}] \phantom{.}
					\begin{compactitem}
						\item Originally brought to my attention by Kaija.
					\end{compactitem}
				\item[{\hyperref[ROMANS-10-13]{Romans 10:13--15, 17}}] \phantom{.}
					\begin{compactitem}
						\item Don't forget verse 17 here!
					\end{compactitem}
				\item[{\hyperref[ROMANS-11-33]{Romans 11:33--36}}]
				\item[{\hyperref[ROMANS-14-13]{Romans 14:13--14}}] \phantom{.}
					\begin{compactitem}
						\item The verses following these also have more stuff about what is ``pure,'' and how it's up to each person.
					\end{compactitem}
				\item[{\hyperref[ROMANS-15-4]{Romans 15:4}}] \phantom{.}
					\begin{compactitem}
						\item On the scriptures specifically.
					\end{compactitem}
				\item[{\hyperref[ROMANS-15-20]{Romans 15:20--21}}]
				\item[{\hyperref[1CORINTHIANS-2]{1 Corinthians 2}}]
				\item[{\hyperref[2CORINTHIANS-3-3]{2 Corinthians 3:3}}]
				\item[{\hyperref[2CORINTHIANS-3-18]{2 Corinthians 3:18}}]
				\item[{\hyperref[2CORINTHIANS-4-3]{2 Corinthians 4:3--7}}]
				\item[{\hyperref[2CORINTHIANS-12-1]{2 Corinthians 12:1--7}}] \phantom{.}
					\begin{compactitem}
						\item This also then leads into a discussion about weakness and strength, which mirrors many points in \hyperref[ETHER-12]{Ether 12}.
					\end{compactitem}
				\item[{\hyperref[2CORINTHIANS-13-1]{2 Corinthians 13:1}}]
				\item[{\hyperref[GALATIANS-1-11]{Galatians 1:11--12}}]
				\item[{\hyperref[GALATIANS-1-15]{Galatians 1:15--16}}]
				\item[{\hyperref[GALATIANS-4-22]{Galatians 4:22--31}}] \phantom{.}
					\begin{compactitem}
						\item Note how Paul is reading a particular story about Abraham as an allegory.
					\end{compactitem}
				\item[{\hyperref[EPHESIANS-1-7]{Ephesians 1:7--9}}]
				\item[{\hyperref[EPHESIANS-1-17]{Ephesians 1:17--19}}]
				\item[{\hyperref[EPHESIANS-3-1]{Ephesians 3:1--12}}]
				\item[{\hyperref[EPHESIANS-3-14]{Ephesians 3:14--21}}] \phantom{.}
					\begin{compactitem}
						\item Doesn't have as much meat as the previous set of verses in this chapter, but it still has some relevant things, including in a connection to an inner transformation.
					\end{compactitem}
				\item[{\hyperref[PHILIPPIANS-1-9]{Philippians 1:9--11}}]
				\item[{\hyperref[COLOSSIANS-1-9]{Colossians 1:9--13}}]
				\item[{\hyperref[COLOSSIANS-1-25]{Colossians 1:25--29}}]
				\item[{\hyperref[COLOSSIANS-2-2]{Colossians 2:2--3}}]
			\end{compactdesc}
			
		% % --
		% \paragraph{Greek New Testament} % *GREEK
		% % --
			% \label{EPIST-REVELATION-GREEK}
		
			% %
			% \subparagraph{\texorpdfstring{\gr{sample word}}{}}
			% %
				% \label{KATALLAGE}
			
				% \begin{compactdesc}
					% \item[BDAG , PDF ] \textbf{}
						% \begin{compactitem}
							% \item 
						% \end{compactitem}
					% \item[LSJ , PDF ] \textbf{}
						% \begin{compactitem}
							% \item 
						% \end{compactitem}
				% \end{compactdesc}
				
				% \begin{compactdesc}
					% \item[Some NT uses] \textbf{}
						% \begin{compactdesc}
							% \item[]
						% \end{compactdesc}
				% \end{compactdesc}
		
	% -----
	\subsubsection{Book of Mormon} % *BOM
	% -----
		\label{EPIST-REVELATION-BOM}
	
		% \begin{tabular}{ll}
			% Total occurrences in the BOM & 
		% \end{tabular}
		
		\begin{compactdesc}
			\item[{\hyperref[1NEPHI-1-1]{1 Nephi 1:1}}] \phantom{.}
				\begin{compactitem}
					\item Nephi here seems to equate the favor he's been shown with the knowledge he's received.
					\item We actually learn quite a bit about where he gets this knowledge, later in the story.
				\end{compactitem}
			\item[{\hyperref[1NEPHI-1-2]{1 Nephi 1:2}}] \phantom{.}
				\begin{compactitem}
					\item See also notes on this verse.
					\item Note the connection to \hyperref[2NEPHI-25-5]{2 Nephi 25:5--6}, where Nephi talks about being taught after the ways of the Jews.
				\end{compactitem}
			\item[{\hyperref[1NEPHI-1-5]{1 Nephi 1:5--7}}] \phantom{.}
				\begin{compactitem}
					\item See also notes on this verse. There are other verses relevant here---it's Lehi's original spiritual manifestation.
					\item It starts with him praying in verse 5.
					\item Note the connection between reading and being filled with the Spirit in verse 12.
					\item His reaction in verse 14 is interesting.
				\end{compactitem}
			\item[{\hyperref[1NEPHI-2-16]{1 Nephi 2:16--19}}] \phantom{.}
				\begin{compactitem}
					\item Compare these verses to \hyperref[DC46-13]{DC 46:13--14}.
					\item Also note in verse 19 the attitude with which Nephi seeks the Lord.
				\end{compactitem}
			\item[{\hyperref[1NEPHI-3-19]{1 Nephi 3:19--20}}] \phantom{.}
				\begin{compactitem}
					\item Certain things needed to be preserved in order for the people to continue to live in a certain way.
					\item Important comments about ``language'' here too.
					\item Verse 20 is also just important for what it says about the spirit and prophets.
				\end{compactitem}
			\item[{\hyperref[1NEPHI-3-29]{1 Nephi 3:29}}] \phantom{.}
				\begin{compactitem}
					\item An important bit for angels.
					\item Also 4:3.
				\end{compactitem}
			\item[{\hyperref[1NEPHI-4-6]{1 Nephi 4:6--38}}] \phantom{.}
				\begin{compactitem}
					\item This whole experience is very important and there's a lot to learn from it. There's a principle in verse 6, and then it plays out in a particular way for Nephi here.
					\item Another thing to pay attention to here is the apparent back and forth Nephi has with the Spirit before killing Laban. But is it an actual back and forth? Is he just remembering this stuff later and inserting the reasoning back into his actions as he reports the story? There's a lot to think about.
				\end{compactitem}
			\item[{\hyperref[1NEPHI-5-4]{1 Nephi 5:4}}]
			\item[{\hyperref[1NEPHI-5-17]{1 Nephi 5:17}}] \phantom{.}
				\begin{compactitem}
					\item Notice that Lehi learns something important, and the Spirit visits him.
				\end{compactitem}
			\item[{\hyperref[1NEPHI-7-14]{1 Nephi 7:14}}]
			\item[{\hyperref[1NEPHI-7-15]{1 Nephi 7:15}}] \phantom{.}
				\begin{compactitem}
					\item A pretty different thing from the previous verse---this is about the ``constrain'' part.
					\item Note that the ``constrain'' language also appears in chapter 4, where Nephi is reasoning about killing Laban.
				\end{compactitem}
			\item[{\hyperref[1NEPHI-10-11]{1 Nephi 10:11}}]
			\item[{\hyperref[1NEPHI-10-17]{1 Nephi 10:17--19}}] \phantom{.}
				\begin{compactitem}
					\item Extremely important verses.
					\item Verse 22 also?
					\item Another thing that's interesting to me here is that Nephi has been listening to his father talk, and wants to know more. And in the next few chapters he does know more, specifically about the vision of the tree. But in chapter 10 here his father has spoken specifically about Christ the Messiah, and about his baptism; Nephi eventually learns about that and records what he learns in 2 Nephi 31. And Lehi here also mentions the olive tree, which Jacob follows up on in Jacob 5. So later chapters in this story are expansions of things that Lehi seems to be the one who was originally talking about them.
					\item In connection with that, note also Nephi's discussion of the olive tree and grafting in chapter 15.
				\end{compactitem}
			\item[{\hyperref[1NEPHI-11-1]{1 Nephi 11:1}}] \phantom{.}
				\begin{compactitem}
					\item Nephi's vision begins in this verse, and continues for a long time. Some good stuff to study.
					\item Especially the way it comes about, and how it happens, with the Spirit speaking to him and everything.
					\item You also have to think about how he went from his experience in chapter 2, to learning more, and eventually having this manifestation, so there's a progression.
					\item And the progression goes with the principles at the end of chapter 10.
					\item 11:14, angel epistemology. Also 11:30.
					\item Note the connection between things like 11:32--33, what Nephi sees, and DC 46. He witnesses it himself. Also 12:6, etc.
					\item 12:18 has some specific stuff about the Holy Ghost.
					\item 13:12--15 has some interesting stuff about the spirit working on people.
					\item 13:29 is important for the scriptures, and also just important for spiritual communication generally. Goes with 34--36, and 40.
					\item 14:1 is about Christ manifesting himself.
					\item Note 14:30 and 15:1 he says he was ``carried away in the spirit.''
				\end{compactitem}
			\item[{\hyperref[1NEPHI-11-11]{1 Nephi 11:11}}] \phantom{.}	
				\begin{compactitem}
					\item This passage is useful for the ``language'' stuff that you get elsewhere, and how God has to adapt the message to our language in order to help us understand.
				\end{compactitem}
			\item[{\hyperref[1NEPHI-13-23]{1 Nephi 13:23--29}}] \phantom{.}
				\begin{compactitem}
					\item This is about the scriptures specifically, but it also talks about the effect that these changes have on people.
					\item It would be interesting to try to figure out whether this view of the scriptures makes sense, given what we know now about the history of the transmission of the manuscripts. I'm not sure how tenable this is. It would make more sense if the taking away of the ``plain and precious things'' had to do with how the scriptures were interpreted, or something like that.
					\item See also verse 32, \hyperref[1NEPHI-14-23]{14:23}, 
				\end{compactitem} 
			\item[{\hyperref[1NEPHI-13-37]{1 Nephi 13:37}}] \phantom{.}
				\begin{compactitem}
					\item A way to receive the Holy Ghost.
				\end{compactitem}
			\item[{\hyperref[1NEPHI-14-18]{1 Nephi 14:18--30}}] \phantom{.}
				\begin{compactitem}
					\item Very interesting remarks about how John would be shown certain things that Nephi also saw, but John would be the one to write them.
					\item I'm especially interested in 25--26, where we learn that other people saw ``all things.'' I assume they would have been shown them in a way that made sense to them, in their language, and so it would be different accounts or stories that would all in reality be about the same thing.
				\end{compactitem}
			\item[{\hyperref[1NEPHI-15-2]{1 Nephi 15:2--4, 6--11}}] \phantom{.}
				\begin{compactitem}
					\item So there's the important part about inquiring of the Lord.
					\item But there's also the fact that Nephi has just had this vision, this enlightenment, and he comes back to people squabbling about little things. Must have seemed very trivial to him, and very pointless.
					\item Make sure to see verses 6--11 here, very important. Verse 11 might be the most important of all.
				\end{compactitem}
			\item[{\hyperref[1NEPHI-15-14]{1 Nephi 15:14}}] \phantom{.}
				\begin{compactitem}
					\item Very important verse for pragmatism and truth and goodness; very important for the principles of revelation as well, for action.
					\item Compare \hyperref[DC93-19]{DC 93:19}.
				\end{compactitem}
			\item[{\hyperref[1NEPHI-15-23]{1 Nephi 15:23--24}}] \phantom{.}
				\begin{compactitem}
					\item About the word of God and the rod.
				\end{compactitem}
			\item[{\hyperref[1NEPHI-16-10]{1 Nephi 16:10}}] \phantom{.}
				\begin{compactitem}
					\item Along with everything connected to the Liahona. This is the verse where it first gets mentioned.
					\item 16:16 is important---this is how we get led step by step.
					\item Verses 27--29 of this chapter are particularly important; this is an extremely important passage to work out for epistemology of revelation stuff.
					\item Notice that the emphasis is on guiding their actions, on getting them to certain places, to perform certain actions, and so on.
					\item See also \hyperref[1NEPHI-18]{1 Nephi 18}, \hyperref[MOSIAH-1-16]{Mosiah 1:16--17} (verse 16 especially is important; it may hint that each person saw something different [``every one''] when they looked at the Liahona---this could be a reason that it was so surprising for them in 1 Nephi 16, when they look at it), 
				\end{compactitem}
			\item[{\hyperref[1NEPHI-17-3]{1 Nephi 17:3}}] \phantom{.}
				\begin{compactitem}
					\item About obedience, but also speaks to the step by step manner in which the Lord will lead us along.
					\item Goes with verse 13.
					\item Verse 30 also.
				\end{compactitem}
			\item[{\hyperref[1NEPHI-17-8]{1 Nephi 17:8}}] \phantom{.}
				\begin{compactitem}
					\item Special knowledge given to Nephi that he wouldn't otherwise be able to have.
				\end{compactitem}
			\item[{\hyperref[1NEPHI-17-13]{1 Nephi 17:13}}]
			\item[{\hyperref[1NEPHI-17-45]{1 Nephi 17:45}}]
			\item[{\hyperref[1NEPHI-18-1]{1 Nephi 18:1--3}}]
			\item[{\hyperref[1NEPHI-18-12]{1 Nephi 18:12--13}}] \phantom{.}
				\begin{compactitem}
					\item And other verses related to this ship journey.
					\item Related to the Liahona too.
				\end{compactitem}
			\item[{\hyperref[1NEPHI-19-12]{1 Nephi 19:12}}] \phantom{.}
				\begin{compactitem}
					\item An important example of context.
				\end{compactitem}
			\item[{\hyperref[1NEPHI-19-20]{1 Nephi 19:20}}] \phantom{.}
				\begin{compactitem}
					\item Some of the physical effects of working with the Spirit.
				\end{compactitem}
			\item[{\hyperref[1NEPHI-19-22]{1 Nephi 19:22--24}}] \phantom{.}
				\begin{compactitem}
					\item Compare \hyperref[2NEPHI-2-14]{2 Nephi 2:14}?
					\item \hyperref[2NEPHI-6-5]{2 Nephi 6:5}
					\item \hyperref[2NEPHI-11-2]{2 Nephi 11:2}, \hyperref[2NEPHI-11-8]{11:8}
				\end{compactitem}
			\item[{\hyperref[1NEPHI-22-1]{1 Nephi 22:1--2}}]
			\item[{\hyperref[2NEPHI-2-4]{2 Nephi 2:4--5}}]
			\item[{\hyperref[2NEPHI-21-2]{2 Nephi 21:2--5}}]
			\item[{\hyperref[2NEPHI-21-9]{2 Nephi 21:9}}]
			\item[{\hyperref[2NEPHI-25-4]{2 Nephi 25:4--7}}] \phantom{.}
				\begin{compactitem}
					\item I'm thinking specifically of the discussion of plainness.
					\item See also verse 20 of this chapter.
				\end{compactitem}
			\item[{\hyperref[2NEPHI-26-1]{2 Nephi 26:1}}]
			\item[{\hyperref[2NEPHI-26-11]{2 Nephi 26:11}}]
			\item[{\hyperref[2NEPHI-26-13]{2 Nephi 26:13}}]
			\item[{\hyperref[2NEPHI-26-22]{2 Nephi 26:22--24}}] \phantom{.}
				\begin{compactitem}
					\item Note verse 33 here as well, related to the ``doing thing in darkness'' point.
				\end{compactitem}
			\item[{\hyperref[2NEPHI-27-23]{2 Nephi 27:23}}] \phantom{.}
				\begin{compactitem}
					\item The ``according to their faith'' part seems crucial. How have I missed this before? The Lord works with us according to our faith, which is basically our level of knowledge. Crucial.
				\end{compactitem}
			\item[{\hyperref[2NEPHI-28]{2 Nephi 28}}] \phantom{.}
				\begin{compactitem}
					\item I read a lot of this chapter as telling us what happens when a person does not have the correct information about how to follow Christ. It goes with passages in restoration scripture that talk about finding the ``true doctrine,'' so that we can know ``how to worship.''
				\end{compactitem}
			\item[{\hyperref[2NEPHI-28-27]{2 Nephi 28:27--31}}]
			\item[{\hyperref[2NEPHI-29-3]{2 Nephi 29:3--14}}] \phantom{.}
				\begin{compactitem}
					\item Not every verse in this block is relevant, but most are, including the important discussion about the Lord giving his word to basically everyone and commanding them to write it.
				\end{compactitem}
			\item[{\hyperref[2NEPHI-30-15]{2 Nephi 30:15--18}}]
			\item[{\hyperref[2NEPHI-31-2]{2 Nephi 31:2--3}}]
			\item[{\hyperref[2NEPHI-31-12]{2 Nephi 31:12}}] \phantom{.}
				\begin{compactitem}
					\item See also \hyperref[2NEPHI-31-17]{31:17}.
				\end{compactitem}
			\item[{\hyperref[2NEPHI-31-19]{2 Nephi 31:19--20}}]
			\item[{\hyperref[2NEPHI-32]{2 Nephi 32}}] \phantom{.}
				\begin{compactitem}
					\item This whole chapter. You can sort of divide it up into three parts: 1--6, 7, and 8--9. All three parts are relevant for the episteomlogy of revelation.
				\end{compactitem}
			\item[{\hyperref[2NEPHI-33-1]{2 Nephi 33:1--2}}]
			\item[{\hyperref[2NEPHI-33-10]{2 Nephi 33:10}}] \phantom{.}
				\begin{compactitem}
					\item See also verse 14; these verses relate to the idea of the words of Christ instructing us, that comes up in chapter 32.
				\end{compactitem}
			\item[{\hyperref[JACOB-1-5]{Jacob 1:5--8}}] \phantom{.}
				\begin{compactitem}
					\item 2:3 mentions ``anxiety'' too.
				\end{compactitem}
			\item[{\hyperref[JACOB-2-5]{Jacob 2:5}}]
			\item[{\hyperref[JACOB-3-2]{Jacob 3:2}}]
			\item[{\hyperref[JACOB-4-5]{Jacob 4:5--18}}] \phantom{.}
				\begin{compactitem}
					\item This is a big block of verses where most of them give some very interesting insights into epist-rev related issues. 
					\item 13 is important for construction and context, and also other things; that verse is very important.
				\end{compactitem}
			\item[{\hyperref[JACOB-6-8]{Jacob 6:8}}]
			\item[{\hyperref[JACOB-7-5]{Jacob 7:5}}] \phantom{.}
				\begin{compactitem}
					\item There are other verses associated with this Sherem story too, like verses 8, 12, 
				\end{compactitem}
			\item[{\hyperref[ENOS-1-2]{Enos 1:2--11}}] \phantom{.}
				\begin{compactitem}
					\item Really Enos's whole experience, especially the initial part.
				\end{compactitem}
			\item[{\hyperref[JAROM-1-4]{Jarom 1:4}}]
			\item[{\hyperref[OMNI-1-20]{Omni 1:20}}]
			\item[{\hyperref[OMNI-1-25]{Omni 1:25}}]
			\item[{\hyperref[WORDSOFMORMON-1-7]{Words of Mormon 1:7--9}}]
			\item[{\hyperref[WORDSOFMORMON-1-16]{Words of Mormon 1:16--18}}] \phantom{.}
				\begin{compactitem}
					\item These verses are about prophets.
				\end{compactitem}
			\item[{\hyperref[MOSIAH-1-3]{Mosiah 1:3--5}}] \phantom{.}
				\begin{compactitem}
					\item This is about the importance of record-keeping.
				\end{compactitem}
			\item[{\hyperref[MOSIAH-2-9]{Mosiah 2:9}}]
			\item[{\hyperref[MOSIAH-2-33]{Mosiah 2:33}}]
			\item[{\hyperref[MOSIAH-2-36]{Mosiah 2:36--37}}]
			\item[{\hyperref[MOSIAH-3-2]{Mosiah 3:2}}] \phantom{.}
				\begin{compactitem}
					\item Here I'm specifically thinking of Benjamin's interaction with this angel that delivers the information.
				\end{compactitem}
			\item[{\hyperref[MOSIAH-3-4]{Mosiah 3:4}}] \phantom{.}
				\begin{compactitem}
					\item Especially the part about Benjamin's righteousness being judged before receiving answers.
				\end{compactitem}
			\item[{\hyperref[MOSIAH-4-9]{Mosiah 4:9}}]
			\item[{\hyperref[MOSIAH-4-11]{Mosiah 4:11--12}}]
			\item[{\hyperref[MOSIAH-5-2]{Mosiah 5:2--4}}]
			\item[{\hyperref[MOSIAH-8-6]{Mosiah 8:6--21}}] \phantom{.}
				\begin{compactitem}
					\item This is about translation and the idea of a ``seer.'' 
				\end{compactitem}
			\item[{\hyperref[MOSIAH-12-25]{Mosiah 12:25--27}}]
			\item[{\hyperref[MOSIAH-13-32]{Mosiah 13:32}}]
			\item[{\hyperref[MOSIAH-18-10]{Mosiah 18:10}}]
			\item[{\hyperref[MOSIAH-18-26]{Mosiah 18:26}}]
			\item[{\hyperref[MOSIAH-21-28]{Mosiah 21:28}}]
			\item[{\hyperref[MOSIAH-26-1]{Mosiah 26:1--3}}]
			\item[{\hyperref[MOSIAH-26-10]{Mosiah 26:10--13}}] \phantom{.}
				\begin{compactitem}
					\item There's more here than this, just wanted to note Alma's experience in being really worried about something and asking the Lord for guidance about it.
				\end{compactitem}
			\item[{\hyperref[MOSIAH-27-10]{Mosiah 27:10--18}}] \phantom{.}
				\begin{compactitem}
					\item This is Alma jr.'s experience with the angel.
					\item Goes with what follows, which is Alma's experience when he's unconscious or whatever.
					\item For this passage, notice in particular that the angel says, I'm here to convince you of the power and authority of God, and not, I'm here to prove to you that God exists. And in verse 18 it seems that this appearance had the desired effect.
				\end{compactitem}
			\item[{\hyperref[ALMA-4-15]{Alma 4:15}}] \phantom{.}
				\begin{compactitem}
					\item The Spirit ``failing'' him because of his own despair or feelings?
					\item Compare to what happens at 8:14--15.
				\end{compactitem}
			\item[{\hyperref[ALMA-4-19]{Alma 4:19--20}}]
			\item[{\hyperref[ALMA-5-45]{Alma 5:45--48}}]
			\item[{\hyperref[ALMA-6-8]{Alma 6:8}}]
			\item[{\hyperref[ALMA-7-4]{Alma 7:4--5}}]
			\item[{\hyperref[ALMA-7-8]{Alma 7:8--9}}]
			\item[{\hyperref[ALMA-7-11]{Alma 7:11--13}}]
			\item[{\hyperref[ALMA-7-14]{Alma 7:14--17}}] \phantom{.}
				\begin{compactitem}
					\item Here I'm thinking specifically of how the Spirit is telling Alma these things he's saying, including in verse 17.
					\item With 18--20, it seems like an example of construction, where Alma is looking around and seeing what the people and the Spirit is giving him revelation connected to what he is perceiving.
					\item See also verse 26.
				\end{compactitem}
			\item[{\hyperref[ALMA-8-10]{Alma 8:10--11}}] \phantom{.}
				\begin{compactitem}
					\item Notice how Alma's prayer and desires interact with what other people choose to do.
				\end{compactitem}
			\item[{\hyperref[ALMA-8-14]{Alma 8:14}}] \phantom{.}
				\begin{compactitem}
					\item Not just this verse, but the following ones also---it's a good example of revelation given for action, without knowing the entire picture, going one step at a time.
					\item Verse 18, for example, all of this is involved.
				\end{compactitem}
			\item[{\hyperref[ALMA-8-20]{Alma 8:20}}] \phantom{.}	
				\begin{compactitem}
					\item And notice how Alma's quest interacts with Amulek here, and the revelations that Amulek has received.
					\item Verses 29--31, etc.
					\item See Alma 10:7--10 for more of Amulek's side of this story, and how he understands his own experiences.
				\end{compactitem}
			\item[{\hyperref[ALMA-8-24]{Alma 8:24}}]
			\item[{\hyperref[ALMA-8-31]{Alma 8:31--32}}] 
			\item[{\hyperref[ALMA-9-19]{Alma 9:19--21}}] \phantom{.}
				\begin{compactitem}
					\item Note the distinction here between the spirit of prophecy and the spirit of revelation, which is expected elsewhere; would be interesting to look into at some point.
					\item Verse 23 as well.
				\end{compactitem} 
			\item[{\hyperref[ALMA-10-5]{Alma 10:5--6}}]
			\item[{\hyperref[ALMA-10-12]{Alma 10:12}}]
			\item[{\hyperref[ALMA-10-17]{Alma 10:17}}]
			\item[{\hyperref[ALMA-10-25]{Alma 10:25}}]
			\item[{\hyperref[ALMA-11-22]{Alma 11:22}}]
			\item[{\hyperref[ALMA-12-3]{Alma 12:3}}] \phantom{.}
				\begin{compactitem}
					\item See verse 7 as well.
				\end{compactitem}
			\item[{\hyperref[ALMA-12-8]{Alma 12:8--11}}] \phantom{.}
				\begin{compactitem}
					\item Extremely important verses.
					\item See my notes on the change in pronouns in verse 10.
					\item Note also that we learn what the ``chains of hell'' are in verse 11: ignorance of God's mysteries.
				\end{compactitem}
			\item[{\hyperref[ALMA-12-28]{Alma 12:28}}] \phantom{.}
				\begin{compactitem}
					\item This verse begins a discussion of how God disseminates knowledge into the world.
					\item See the verses following this for more. It extends into chapter 13.
						\begin{compactitem}
							\item You see it again in 13:22 and 13:24, for example, and 13:26.
						\end{compactitem}
					\item A similar discussion is in \hyperref[MORONI-7]{Moroni 7}. The more I think about this stuff, the more important I realize this is.
				\end{compactitem}
			\item[{\hyperref[ALMA-12-30]{Alma 12:30}}]
			\item[{\hyperref[ALMA-13-3]{Alma 13:3--10}}] \phantom{.}	
				\begin{compactitem}
					\item Still connected to the issues above about disseminating knowledge of the plan of salvation.
				\end{compactitem}
			\item[{\hyperref[ALMA-13-28]{Alma 13:28}}]
			\item[{\hyperref[ALMA-14-10]{Alma 14:10--11}}]
			\item[{\hyperref[ALMA-16-4]{Alma 16:4--6}}]
			\item[{\hyperref[ALMA-16-15]{Alma 16:15--17}}]
			\item[{\hyperref[ALMA-17-2]{Alma 17:2--4}}] \phantom{.}
				\begin{compactitem}
					\item Related to this, see verses 9--10 of the same chapteFr, and then the Lord begins speaking to them in verse 11.
				\end{compactitem}
			\item[{\hyperref[ALMA-18-2]{Alma 18:2}}] \phantom{.}
				\begin{compactitem}
					\item Here begins Ammon's preaching to Lamoni. I think there are some examples of construction here, where the preaching and revelation occurs on the basis of what Lamoni already believes.
					\item Other verses relevant here are 18, 24--32, 36ff, 
					\item And notice how this interacts with the ideas on the ``traditions of the fathers'' in DC 93.
				\end{compactitem}
			\item[{\hyperref[ALMA-18-16]{Alma 18:16}}] \phantom{.}
				\begin{compactitem}
					\item This is where Ammon begins to receive insight from the Holy Ghost in his preaching to Lamoni.
					\item Other verses that are relevant here include 18, 20,
				\end{compactitem}
			\item[{\hyperref[ALMA-18-35]{Alma 18:35}}] \phantom{.}
				\begin{compactitem}
					\item There are a few things to note here. 
					\item One is that a \textit{portion} of the Spirit dwells with him, and so we can have more or less of it. 
					\item Another is that the way the Spirit is described here is similar to how the Liahona is described, which is important when you think of externalized sources of revelation, like the Liahona and seer stones.
				\end{compactitem}
			\item[{\hyperref[ALMA-19-6]{Alma 19:6}}]
			\item[{\hyperref[ALMA-19-16]{Alma 19:16--17}}] \phantom{.}
				\begin{compactitem}
					\item Notice Abish here acting as an ``angel'' or minister to share about Christ.
				\end{compactitem}
			\item[{\hyperref[ALMA-20-4]{Alma 20:4--6}}]
			\item[{\hyperref[ALMA-21-16]{Alma 21:16--17}}]
			\item[{\hyperref[ALMA-22]{Alma 22}}] \phantom{.}
				\begin{compactitem}
					\item Here begins the preaching of Aaron to Lamoni's father.
					\item It involves the Spirit in various ways, including in Aaron being led there and also in the guy's questions.
					\item It's also a good example of construction, or building upon what is already there. Notice that Aaron doesn't correct the king's views, he just builds on them.
				\end{compactitem}
			\item[{\hyperref[ALMA-23-6]{Alma 23:6}}]
			\item[{\hyperref[ALMA-24-7]{Alma 24:7--9}}] \phantom{.}
				\begin{compactitem}
					\item I think I have notes on this elsewhere, but all the scriptures involving the ``traditions of the fathers'' are relevant for the epistemology of revelation; see DC 93, and also think about the maze of many starting-points.
				\end{compactitem}
			\item[{\hyperref[ALMA-24-14]{Alma 24:14}}] \phantom{.}
				\begin{compactitem}
					\item There is a reference here to the ``angels'' that make known the plan of salvation.
					\item Two other places that talk about this a lot are Alma 12 and Moroni 7.
					\item See also \hyperref[ALMA-27-4]{Alma 27:4}.
					\item This is a major theme in different places of the Book of Mormon and I should write about it eventually. It's like a big-picture view of how knowledge of Christ gets into the world in the first place.
				\end{compactitem}
			\item[{\hyperref[ALMA-24-29]{Alma 24:29--30}}]
			\item[{\hyperref[ALMA-25-15]{Alma 25:15--16}}] \phantom{.}
				\begin{compactitem}
					\item Important for construction and accountability.
				\end{compactitem}
			\item[{\hyperref[ALMA-26-14]{Alma 26:14--15}}] \phantom{.}
				\begin{compactitem}
					\item Note the ``chains of hell'' here and how it has to do with ignorance.
				\end{compactitem}
			\item[{\hyperref[ALMA-26-21]{Alma 26:21--22}}]
			\item[{\hyperref[ALMA-27-7]{Alma 27:7--13}}] \phantom{.}
				\begin{compactitem}
					\item Another experience of someone asking the Lord something on behalf of a bunch of other people, and the Lord responding in a particular way.
				\end{compactitem}
			\item[{\hyperref[ALMA-28-12]{Alma 28:12}}]
			\item[{\hyperref[ALMA-29-4]{Alma 29:4--8}}] \phantom{.}
				\begin{compactitem}
					\item There's kind of two blocks of verses here: 4--5, and 6--8. 8 in particular is important for the dissemination of gospel knowledge at population levels.
				\end{compactitem} 
			\item[{\hyperref[ALMA-30-13]{Alma 30:13--60}}] \phantom{.}
				\begin{compactitem}
					\item Here begins Korihor's preaching, which has a lot to do with what can be known or not.
					\item There are many interesting points here, and Alma and Amulek pick up on so many of these issues in the next few chapters.
				\end{compactitem}
			\item[{\hyperref[ALMA-31-5]{Alma 31:5}}]
			\item[{\hyperref[ALMA-32]{Alma 32}}]
				\begin{compactitem}
					\item An extremely important chapter; one of the most important of all.
					\item Compare the ``sower sowing the word'' parable from Christ.
				\end{compactitem}
			\item[{\hyperref[ALMA-33-19]{Alma 33:19--21}}]
			\item[{\hyperref[ALMA-34-33]{Alma 34:33--39}}]
			\item[{\hyperref[ALMA-36-4]{Alma 36:4--5}}]
			\item[{\hyperref[ALMA-36-18]{Alma 36:18--26}}] \phantom{.}
				\begin{compactitem}
					\item Could also be relevant for Alma 32.
				\end{compactitem}
			\item[{\hyperref[ALMA-37-4]{Alma 37:4--5}}] \phantom{.}
				\begin{compactitem}
					\item The ``brightness'' of the records matters; something that is bright gives off light.
				\end{compactitem}
			\item[{\hyperref[ALMA-37-6]{Alma 37:6--12}}] \phantom{.}
				\begin{compactitem}
					\item Lots about the scriptures, mysteries, and other things. Very rich passages.
					\item And repentance as learning, clearly stated in verse 9. There is a strong connection between repentance and learning or gaining knowledge.
				\end{compactitem}
			\item[{\hyperref[ALMA-37-13]{Alma 37:13--37}}] \phantom{.}
				\begin{compactitem}
					\item A long passage with very interesting details about how Helaman is instructed to disseminate knowledge relevant to the gospel.
					\item This is important for understanding things like Alma 12:9--11 and other passages.
				\end{compactitem}
			\item[{\hyperref[ALMA-37-38]{Alma 37:38--47}}]
			\item[{\hyperref[ALMA-38-6]{Alma 38:6--10}}] \phantom{.}
				\begin{compactitem}
					\item Verse 6 is probably most important here.
				\end{compactitem}
			\item[{\hyperref[ALMA-39-5]{Alma 39:5--6}}]
			\item[{\hyperref[ALMA-39-15]{Alma 39:15--19}}] \phantom{.}
				\begin{compactitem}
					\item More about dissemination, priesthood callings to share knowledge, etc. I'm noticing these patterns and need to explore them at some point.
					\item Here you can connect to the ``angels'' stuff through the word ``ministry'' in 16.
					\item And angels are mentioned in 19 too! Great passage.
				\end{compactitem}
			\item[{\hyperref[ALMA-40-3]{Alma 40:3}}] \phantom{.}
				\begin{compactitem}
					\item Make sure to see his follow-up to this in verses 9 and 11.
				\end{compactitem}
			\item[{\hyperref[ALMA-40-13]{Alma 40:13}}]
			\item[{\hyperref[ALMA-43-2]{Alma 43:2}}]
			\item[{\hyperref[ALMA-43-23]{Alma 43:23--24}}]
			\item[{\hyperref[ALMA-45-9]{Alma 45:9--15}}] \phantom{.}
				\begin{compactitem}
					\item Alma's prophecy is important here, but so is the fact that this looks like an instance of the principle in Alma 12:9--11, where certain people are commanded not to reveal particular things they know.
					\item 12 is also important, going against knowledge.
				\end{compactitem} 
			\item[{\hyperref[ALMA-46-8]{Alma 46:8--9}}]
			\item[{\hyperref[ALMA-47-36]{Alma 47:36}}] \phantom{.}
				\begin{compactitem}
					\item See also 48:24 for a related idea.
				\end{compactitem}
			\item[{\hyperref[ALMA-48-11]{Alma 48:11}}]
			\item[{\hyperref[ALMA-49-30]{Alma 49:30}}]
			\item[{\hyperref[ALMA-55-1]{Alma 55:1}}]
			\item[{\hyperref[ALMA-57-21]{Alma 57:21}}]
			\item[{\hyperref[HELAMAN-3-29]{Helaman 3:29}}] \phantom{.}
				\begin{compactitem}
					\item See notes on this verse and the connection to Hebrews; very important for the maze of many starting-points and other issues.
				\end{compactitem} 
			\item[{\hyperref[HELAMAN-3-35]{Helaman 3:35}}] \phantom{.}
				\begin{compactitem}
					\item Especially with the ``faith'' part.
				\end{compactitem}
			\item[{\hyperref[HELAMAN-4-12]{Helaman 4:12--13}}]
			\item[{\hyperref[HELAMAN-4-20]{Helaman 4:20--26}}]
			\item[{\hyperref[HELAMAN-5-17]{Helaman 5:17--18}}]
			\item[{\hyperref[HELAMAN-5-29]{Helaman 5:29}}] \phantom{.}
				\begin{compactitem}
					\item And verses after this one---the voice that comes, its nature and what it says.
					\item And verse 33, about words that can't be uttered by a person.
					\item This entire story is full of interesting points to consider.
					\item 41 is important, and seems to go with the ``particle of faith'' remark in Alma 32. I wonder if you can impose that Alma 32 model on what happens to them in this chapter.
					\item This is also a helpful place to look at the role of angels in giving knowledge.
					\item Verses 49--50---to some it is given to know, to others it is given to believe on their words...
				\end{compactitem}
			\item[{\hyperref[HELAMAN-6-34]{Helaman 6:34--36}}]
			\item[{\hyperref[HELAMAN-7-23]{Helaman 7:23--24}}]
			\item[{\hyperref[HELAMAN-7-29]{Helaman 7:23--29}}]
			\item[{\hyperref[HELAMAN-8-14]{Helaman 8:14--25}}] \phantom{.}
				\begin{compactitem}
					\item Many different points made about many different issues here.
				\end{compactitem}
			\item[{\hyperref[HELAMAN-9-1]{Helaman 9:1--5}}]
			\item[{\hyperref[HELAMAN-11-23]{Helaman 11:23}}]
			\item[{\hyperref[HELAMAN-12]{Helaman 12}}] \phantom{.}
				\begin{compactitem}
					\item There's a large stretch of verses in the middle of this chapter, about God commanding things to happen, that aren't relevant for revelation or epistemology. But many of the other parts are. We basically get to decide what role we want God to play in our life, and this affects what knowledge we can receive.
				\end{compactitem}
			\item[{\hyperref[HELAMAN-13-3]{Helaman 13:3--5}}] 
			\item[{\hyperref[HELAMAN-13-7]{Helaman 13:7--8}}] \phantom{.}
				\begin{compactitem}
					\item Note that these verses are here for different reasons---7 has to do with the angel/dissemination thing.
					\item In 8, notice that losing the word and the Spirit are two separate things.
					\item This is a good example of the angel epistemology thing, which I need to take up someday.
				\end{compactitem}
			\item[{\hyperref[HELAMAN-14-28]{Helaman 14:28--31}}]
			\item[{\hyperref[HELAMAN-15-4]{Helaman 15:4--17}}] \phantom{.}
				\begin{compactitem}
					\item The idea of the ``true knowledge'' is important---it seems to refer specifically to Christ. It's the knowledge you need in order to worship correctly and progress.
				\end{compactitem} 
			\item[{\hyperref[HELAMAN-16-1]{Helaman 16:1--5}}] \phantom{.}
				\begin{compactitem}
					\item Notice here the various things that convince people. What matters is just being convinced and acting---getting baptized---and not anything else.
				\end{compactitem}
			\item[{\hyperref[HELAMAN-16-14]{Helaman 16:14}}] \phantom{.}
				\begin{compactitem}
					\item More angel epistemology.
				\end{compactitem}
			\item[{\hyperref[HELAMAN-16-15]{Helaman 16:15--23}}] \phantom{.}
				\begin{compactitem}
					\item I think what's important about this story is that you can kind of convince yourself of anything, in any direction, no matter what the evidence is; this is as true for religious matters as it is for the history of science or politics or anything else. Often what drives our reasoning is not the evidence but our attitude toward where the evidence might lead.
					\item This is a feature of human reasoning, not necessarily a bug; I have to think more about what it means for religious matters.
				\end{compactitem}
			\item[{\hyperref[3NEPHI-1-22]{3 Nephi 1:22}}] \phantom{.}
				\begin{compactitem}
					\item I'm thinking specifically here of how the people react to the sign.
					\item This also goes with 3 Nephi 2:1--3 below.
				\end{compactitem}
			\item[{\hyperref[3NEPHI-1-24]{3 Nephi 1:24--25}}]
			\item[{\hyperref[3NEPHI-1-29]{3 Nephi 1:29--30}}]
			\item[{\hyperref[3NEPHI-2-1]{3 Nephi 2:1--3}}]
			\item[{\hyperref[3NEPHI-3-19]{3 Nephi 3:19}}]
			\item[{\hyperref[3NEPHI-5-20]{3 Nephi 5:20}}] \phantom{.}
				\begin{compactitem}
					\item Knowledge is given for salvation.
				\end{compactitem}
			\item[{\hyperref[3NEPHI-6-18]{3 Nephi 6:18}}] \phantom{.}
				\begin{compactitem}
					\item In other words, they knew what to \textit{do}, which is the key knowledge.
				\end{compactitem}
			\item[{\hyperref[3NEPHI-6-20]{3 Nephi 6:20}}] \phantom{.}
				\begin{compactitem}
					\item Angels/dissemination stuff.
				\end{compactitem}
			\item[{\hyperref[3NEPHI-7-15]{3 Nephi 7:15--16}}] \phantom{.}
				\begin{compactitem}
					\item Very important for the seeing stuff---he sees the ministry of Christ, or something like that. In some way he knows about it.
					\item Then he goes and preaches.
				\end{compactitem}
			\item[{\hyperref[3NEPHI-7-18]{3 Nephi 7:18--22}}] \phantom{.}
				\begin{compactitem}
					\item Also important for seeing, since this is the result of Nephi's preaching.
					\item This pattern is repeated throughout the Book of Mormon, and at some point I'd like to study it in detail.
				\end{compactitem}
			\item[{\hyperref[3NEPHI-9-1]{3 Nephi 9:1}}] \phantom{.}
				\begin{compactitem}
					\item Here the people begin to hear the voice of Christ.
				\end{compactitem}
			\item[{\hyperref[3NEPHI-10-18]{3 Nephi 10:18--19}}]
			\item[{\hyperref[3NEPHI-11]{3 Nephi 11}}] \phantom{.}
				\begin{compactitem}
					\item There are many important things here, including the voice, the number of times they hear it, the fact that Christ manifests himself to the people, the fact that they see him and touch him, their certainty, the fact that the text says that they will know by doing so, and so on. Lots of important things in this chapter.
					\item Including 31--40 in issues related to doctrine and receiving it.
				\end{compactitem}
			\item[{\hyperref[3NEPHI-12-1]{3 Nephi 12:1--2}}] \phantom{.}
				\begin{compactitem}
					\item More ``see and know.''
					\item Notice also the connection here to DC 46.
				\end{compactitem}
			\item[{\hyperref[3NEPHI-14-7]{3 Nephi 14:7--8}}]
			\item[{\hyperref[3NEPHI-14-24]{3 Nephi 14:24--}}] \phantom{.}
				\begin{compactitem}
					\item On the importance of action.
				\end{compactitem}
			\item[{\hyperref[3NEPHI-15-18]{3 Nephi 15:18}}]
			\item[{\hyperref[3NEPHI-15-22]{3 Nephi 15:22--24}}]
			\item[{\hyperref[3NEPHI-16-1]{3 Nephi 16:1--4}}]
			\item[{\hyperref[3NEPHI-16-6]{3 Nephi 16:6--7}}]
			\item[{\hyperref[3NEPHI-16-10]{3 Nephi 16:10--12}}]
			\item[{\hyperref[3NEPHI-17-2]{3 Nephi 17:2--3}}]
			\item[{\hyperref[3NEPHI-17-5]{3 Nephi 17:5--8}}] \phantom{.}
				\begin{compactitem}
					\item Notice how Christ responds to what he sees the people wanting.
				\end{compactitem}
			\item[{\hyperref[3NEPHI-17-15]{3 Nephi 17:15--18}}]
			\item[{\hyperref[3NEPHI-17-24]{3 Nephi 17:24--25}}]
			\item[{\hyperref[3NEPHI-18-7]{3 Nephi 18:7}}] \phantom{.}
				\begin{compactitem}
					\item Also verse 11, which says something similar.
				\end{compactitem}
			\item[{\hyperref[3NEPHI-18-25]{3 Nephi 18:25}}]
			\item[{\hyperref[3NEPHI-19-2]{3 Nephi 19:2--3}}] \phantom{.}
				\begin{compactitem}
					\item Another example of DC 46 stuff, about people knowing and then sharing the knowledge.
					\item Verse 21 too, and possibly others. Verse 21 is actually a very important example of it.
					\item 28 also.
				\end{compactitem}
			\item[{\hyperref[3NEPHI-19-9]{3 Nephi 19:9--36}}] \phantom{.}
				\begin{compactitem}
					\item This is a whole episode that is worth looking at closely.
				\end{compactitem}
			\item[{\hyperref[3NEPHI-20-8]{3 Nephi 20:8--9}}]
			\item[{\hyperref[3NEPHI-20-27]{3 Nephi 20:27}}]
			\item[{\hyperref[3NEPHI-21-2]{3 Nephi 21:2}}]
			\item[{\hyperref[3NEPHI-21-6]{3 Nephi 21:6}}]
			\item[{\hyperref[3NEPHI-26-6]{3 Nephi 26:6--12}}] \phantom{.}
				\begin{compactitem}
					\item Very important application of other principles.
					\item And when we will know these things, if we're faithful? I think the ansewr is, when Christ himself teaches them to us, either in person or through the Holy Ghost.
				\end{compactitem}
			\item[{\hyperref[3NEPHI-26-14]{3 Nephi 26:14}}] \phantom{.}
				\begin{compactitem}
					\item Also verses 16, 18
				\end{compactitem}
			\item[{\hyperref[3NEPHI-27-1]{3 Nephi 27:1--2}}] \phantom{.}
			\item[{\hyperref[3NEPHI-28-11]{3 Nephi 28:11}}]
			\item[{\hyperref[3NEPHI-28-12]{3 Nephi 28:12--16}}]
			\item[{\hyperref[3NEPHI-28-29]{3 Nephi 28:29}}]
			\item[{\hyperref[3NEPHI-28-30]{3 Nephi 28:30}}]
			\item[{\hyperref[3NEPHI-28-34]{3 Nephi 28:34--35}}]
			\item[{\hyperref[3NEPHI-29-6]{3 Nephi 29:6}}]
			\item[{\hyperref[4NEPHI-1-38]{4 Nephi 1:38--39}}]
			\item[{\hyperref[MORMON-1-14]{Mormon 1:14--17}}] \phantom{.}
				\begin{compactitem}
					\item Several different important things in this block.
				\end{compactitem}
			\item[{\hyperref[MORMON-2-26]{Mormon 2:26}}]
			\item[{\hyperref[MORMON-3-16]{Mormon 3:16}}]
			\item[{\hyperref[MORMON-3-20]{Mormon 3:20--21}}]
			\item[{\hyperref[MORMON-4-11]{Mormon 4:11}}] \phantom{.}
				\begin{compactitem}
					\item Compare this to other descriptions of being not being able to say something; this one seems different.
				\end{compactitem}
			\item[{\hyperref[MORMON-5-8]{Mormon 5:8--9}}]
			\item[{\hyperref[MORMON-5-12]{Mormon 5:12--15}}] \phantom{.}
				\begin{compactitem}
					\item This is about the purpose of the scriptures.
				\end{compactitem}
			\item[{\hyperref[MORMON-5-16]{Mormon 5:16--18}}] \phantom{.}
				\begin{compactitem}
					\item Note particularly in verse 17, that it talks about being led by God the Father; see \hyperref[DC84-48]{DC 84:48}.
				\end{compactitem}
			\item[{\hyperref[MORMON-8-12]{Mormon 8:12}}]
			\item[{\hyperref[MORMON-8-14]{Mormon 8:14--17}}]
			\item[{\hyperref[MORMON-8-34]{Mormon 8:34--36}}]
			\item[{\hyperref[MORMON-9-7]{Mormon 9:7--9}}]
			\item[{\hyperref[MORMON-9-21]{Mormon 9:21}}]
			\item[{\hyperref[MORMON-9-25]{Mormon 9:25--}}] \phantom{.}
				\begin{compactitem}
					\item Very important, especially 25; I think the rest explains how the confirmation of his words is to take place. Goes with the kind of Schiller stuff in ``Axioms as Postulates,'' PDF 18.
				\end{compactitem}
			\item[{\hyperref[ETHER-2]{Ether 2--3}}] (See the \hyperref[TALKIDEAS]{talk idea} about the preparing of the stones as a metaphor.)
				\begin{compactitem}
					\item There are also a lot of other things going on this these two chapters, that have to do with specific ways the Lord is interacting with people and instructing them.
					\item I've tried to get at some of them with the scriptures listed below.
					\item 3:4 starts when the Lord puts his hand out. This and the rest of the chapter are extremely important.
				\end{compactitem}
			\item[{\hyperref[ETHER-2-12]{Ether 2:12}}] \phantom{.}
				\begin{compactitem}
					\item One of the ways that Jesus Christ is manifest to us is through the writings of others. This would be important for thinking about other places where it talks about him being ``manifest,'' or related things being manifest; I think there are places in 2 Nephi 25 or 26, or around there.
				\end{compactitem}
			\item[{\hyperref[ETHER-2-13]{Ether 2:13--15}}]
			\item[{\hyperref[ETHER-2-25]{Ether 2:25--3:}}] \phantom{.}
				\begin{compactitem}
					\item This is the specific story of the stones, with the Lord asking him what he wants to do, and Jared's brother coming up with a plan where he makes the stones, and so on.
				\end{compactitem}
			\item[{\hyperref[ETHER-4-1]{Ether 4:1--3}}] \phantom{.}
				\begin{compactitem}
					\item Regarding how the records are going to be shared.
				\end{compactitem}
			\item[{\hyperref[ETHER-4]{Ether 4--5}}] \phantom{.}
				\begin{compactitem}
					\item Taken together with what happens in chapter 3, Ether 4 is one of the most important chapters about epistemological stuff. There are so many important points here.
					\item Chapter 5 is short but has some good stuff too, but I'm just wowed by chapter 4.
				\end{compactitem}
			\item[{\hyperref[ETHER-6-1]{Ether 6:1--12}}] \phantom{.}
				\begin{compactitem}
					\item It's a nice story of being pushed in a general direction, even if it seems like it's chaotic.
				\end{compactitem}
			\item[{\hyperref[ETHER-9-22]{Ether 9:22}}]
			\item[{\hyperref[ETHER-11-13]{Ether 11:13}}] \phantom{.}
				\begin{compactitem}
					\item The pattern of prophets is here, and elsewhere in Ether, you just have to search for ``prophet'' in the book.
				\end{compactitem}
			\item[{\hyperref[ETHER-12-3]{Ether 12:3--41}}] \phantom{.}
				\begin{compactitem}
					\item On faith, but I think there's some very important things here about the principle of finding what you look for.
					\item Verse 35 is sneakily important, I think.
					\item Verse 39, for language stuff.
				\end{compactitem}
			\item[{\hyperref[ETHER-15-19]{Ether 15:19}}]
			\item[{\hyperref[MORONI-2-2]{Moroni 2:2--3}}]
			\item[{\hyperref[MORONI-7-4]{Moroni 7:4--19}}] \phantom{.}
				\begin{compactitem}
					\item This block of verses is all about discerning good from evil in order to allow us to act, and in particular, to act to ``lay hold on every good thing.''
					\item So the whole first part of the chapter is about the way we make that particular discernment.
				\end{compactitem}
			\item[{\hyperref[MORONI-7-22]{Moroni 7:22--39}}] \phantom{.}
				\begin{compactitem}
					\item Very important passage for the issue of angels and dissemination of knowledge.
					\item 23 also includes the difference with prophets.
					\item 25 is crucial.
					\item This is so important; I just keep thinking about how these brief comments about it are made throughout the Book of Mormon, and I really need to tie them together at some point.
					\item It's also very important to get the parts in here about faith and how it's an active force that makes things happen and shapes what is possible for us to do in the world; verse 38 is a good example of this.
				\end{compactitem}
			\item[{\hyperref[MORONI-7-44]{Moroni 7:44}}]
			\item[{\hyperref[MORONI-8-7]{Moroni 8:7--9}}]
			\item[{\hyperref[MORONI-8-22]{Moroni 8:22}}] \phantom{.}
				\begin{compactitem}
					\item This is absolutely fundamental to what it means for the law to exist for someone; you have to KNOW ABOUT IT in order to be held accountable to it. This goes hand in hand with the restoration definition of truth in terms of knowledge.
					\item Verse 24 also, equally important.
				\end{compactitem}
			\item[{\hyperref[MORONI-8-25]{Moroni 8:25--26}}]
			\item[{\hyperref[MORONI-8-28]{Moroni 8:28--29}}]
			\item[{\hyperref[MORONI-9-4]{Moroni 9:4}}]
			\item[{\hyperref[MORONI-9-20]{Moroni 9:20}}]
			\item[{\hyperref[MORONI-10-3]{Moroni 10:3--7}}]
			\item[{\hyperref[MORONI-10-8]{Moroni 10:8--19}}] \phantom{.}
				\begin{compactitem}
					\item On spiritual gifts.
				\end{compactitem}
		\end{compactdesc}
		
	% -----
	\subsubsection{Doctrine and Covenants} % *DC
	% -----
		\label{EPIST-REVELATION-DC}
	
		% \begin{tabular}{ll}
			% Total occurrences in the DC & 
		% \end{tabular}
		
		\begin{compactdesc}
			\item \hyperref[DC1-24]{DC 1:24}
				\begin{compactitem}
					\item See related verses in the BOM about ``weakness,'' and how what they say connects to the revelatory process we see with JS.
				\end{compactitem}
			\item \hyperref[DC1-25]{DC 1:25--26, 28}
				\begin{compactitem}
					\item Humility given here as a key to receive knowledge.
					\item Notice especially the ``receive knowledge from time to time'' bit at the end of 28.
				\end{compactitem}
			\item \hyperref[DC1-33]{DC 1:33}
			\item \hyperref[DC3-4]{DC 3:4--15}
				\begin{compactitem}
					\item See also the JSP version of some verses, including 12.
					\item This whole story is a powerful example of how setting aside righteous behavior turns off the fountain of revelation in our lives.
				\end{compactitem}
			\item \hyperref[DC5]{DC 5} 
				\begin{compactitem}
					\item The issue of seeing the plates is extremely important, along with the ``power'' involved; this deserves a detailed treatment at some point in the future.
					\item It comes up again in verses 24--26 and 28 as well.
				\end{compactitem}
			\item \hyperref[DC5-11]{DC 5:11--14}
				\begin{compactitem}
					\item Note the JSP manuscript.
					\item Other verses in this section are relevant too, like 24--26; this entire experience of Martin Harris is a good lesson on revelation and some principles involved.
				\end{compactitem}
			\item \hyperref[DC5-16]{DC 5:16}
				\begin{compactitem}
					\item Interesting little principle here about belief and visitation. Seems like it applies to everyone.
				\end{compactitem}
			\item \hyperref[DC6-7]{DC 6:7--1}
				\begin{compactitem}
					\item A potentially very powerful principle.
					\item Verse 11 also mentions the mysteries.
					\item You have to seek wisdom, and you have to inquire; otherwise you aren't going to know them. You have to look out for them and want to know them and be seeking answers.
				\end{compactitem}
			\item \hyperref[DC6-14]{DC 6:14--17}
				\begin{compactitem}
					\item Need also to think of these verses in the context of Cowdery's life at the time.
					\item With verse 15, Bednar said some things about it in the training that Brett showed us in December 2021 at a bishopric meeting. 
					\item Verse 17, the Lord has shared these things with him for a specific purpose: to covince him that the words he's been writing, as part of the BOM, are true.
				\end{compactitem}
			\item \hyperref[DC6-22]{DC 6:22--24}
			\item \hyperref[DC6-22]{DC 6:25--28}
				\begin{compactitem}
					\item On translation, but note that verse 26 has a principle about what people receive in general, like Alma 12.
				\end{compactitem}
			\item \hyperref[DC6-28]{DC 6:28}
				\begin{compactitem}
					\item I am not totally sure what to make of the ``two or three witnesses'' principle. It might relate to social epistemology and sharing beliefs with others, rather than establishing the actual truth of something.
				\end{compactitem}
			\item \hyperref[DC7]{DC 7}
				\begin{compactitem}
					\item The changes over time to this section, from the JSP/BC version to now, are pretty interesting, and clearly reflect a growth of knowledge and understanding of other gospel principles that are then inserted later, into the 1835 DC version.
					\item Also, the weird thing about the ``parchment'' and the translation, which is not mentioned in the earliest versions.
				\end{compactitem}
			\item \hyperref[DC8]{DC 8}
				\begin{compactitem}
					\item First, look at the note at verse 6 on the use of ``sprout,'' ``rod,'' and ``gift of Aaron'' in successive versions---it's possible to link the modern version of the verse into stuff about the priesthood and the ministering of angels (and therefore the epistemology of revelation).
					\item Another thing to look at here is this verse (6) itself and its evolution over time---could be another example of how JS's (and everyone else's) evolving understanding of the gospel led them to revise a revelation in order to better explain and correspond to that evolving understanding.
					\item And yet another thing is the fact that these people were using things like rods in order to get knowledge. The Lord is using what is there to instruct.
						\begin{compactitem}
							\item Verse 7, and its original (about ``Nature''), is important here.
							\item Verse 8 also, especially the original---what is the antecedent of ``it''? Seems to be referring to the ``thing of Nature'' working in Cowdery's hands.
						\end{compactitem}
					\item Verses 1--3 also have important principles about revelation that go beyond this section too. Connecting to the ``spirit of revelation,'' and the connection especially to Moses and his experience, is really interesting and there would be a lot to look at there.
					\item The later verses, 9--11, also have principles about revelation, including Cowdery being specifically instructed to ask in order to ``know the mysteries.''
					\item Verse 1 at the end, I think it's talking about what it says about the records in the BOM, which would make the BOM the result of ``the manifestation of my spirit.''
				\end{compactitem}
			\item \hyperref[DC9]{DC 9}
				\begin{compactitem}
					\item Verse 5, continuing as you commenced---this is probably a more general principle that isn't unique to the process of receiving revelation. Though I think it's relevant for revelation too.
					\item Verse 7, with my note.
					\item This is one of the most important passages of scripture anywhere for the epistemology of revelation.
				\end{compactitem}
			\item \hyperref[DC10]{DC 10}
				\begin{compactitem}
					\item Note especially the JSP background to this revelation---how it shares language in common with certain passages from the BOM, and it isn't clear when it was received. There is a construction or production process for this revelation that takes place over time and has many different influences.
						\begin{compactitem}
							\item There are other aspects of this to note too, for example, the inclusion of a reference to the ``Urim and Thummim'' that was probably not there originally (it isn't there in the BC).
							\item There are other important differences too.
						\end{compactitem}
					\item The idea of re-translating is also interesting, when you think of translation as governed by the rules of revelation.
					\item Verse 45---some parts of scripture give you ``greater views'' of the gospel.
					\item Verse 46, an application of the Alma 12 principle.
						\begin{compactitem}
							\item Notice that the way you understand ``this work,'' as either the BOM or the work of restoration more generally, tells you something about the need for more revelations after the BOM. The BOM was the gospel that prior prophets wanted people to have, but there was more to have than that.
						\end{compactitem}
					\item Verse 52, clear example of construction.
				\end{compactitem}
			\item DC 11
				\begin{compactitem}
					\item \hyperref[DC11-12]{DC 11:12--14}
					\item \hyperref[DC11-21]{DC 11:21}
						\begin{compactitem}
							\item Notice \hyperref[DC23-3]{DC 23:3}, where this promise appears to be fulfilled (or at least it's followed up on).
						\end{compactitem}
					\item \hyperref[DC11-22]{DC 11:22} This verse is very important as a complement to \hyperref[ALMA-12]{Alma 12}, showing that there is a general amount of the word (or light) given to the world at large. 
						\begin{compactitem}
							\item It's also important because it's an example of construction, or it's related to that. All things will be added on top of what's already been given---but if what's already been given depended on what people knew before, and what's coming after has to be in line with what came before, then what came before has a shaping effect on what comes after. That's part of construction.
							\item Verse 21 tells Hyrum to obtain the word, and talks about the resulting desires.
							\item Verse 22 then says how to obtain the word: you have to study, you have to fill your mind with real material.
						\end{compactitem}
				\end{compactitem}
			\item DC 18
				\begin{compactitem}
					\item \hyperref[DC18-1]{DC 18:1--4}
					\item \hyperref[DC18-18]{DC 18:18}
						\begin{compactitem}
							\item EXTREMELY important verse about the Spirit manifesting things which are ``expedient'' to us.
						\end{compactitem}
					\item \hyperref[DC18-33]{DC 18:33--36}
						\begin{compactitem}
							\item This is an interesting discussion of the voice of Christ, and its relation to the Spirit. Verse 35 in particular seems to suggest that, because the revelations have been given by the voice of Christ, the people who read them can then testify (verse 36) that they have heard the voice of Christ. Pretty interesting. Related to the issue of DC 46.
							\item Verse 35 also says that Christ's voice is the one speaking, but the words are given by the Spirit. So that counts as hearing Christ's voice.
							\item Also makes me think of \hyperref[2NEPHI-32]{2 Nephi 32}, especially \hyperref[2NEPHI-32-6]{verse 6}, and how to take the bit at the end of that verse about the things which Christ will ``say'' unto us. If we go with this way of understanding these verses in DC 18, then you can see us as receiving the revelations and through them hearing the voice of Christ and finding out what he is saying to us. See also \hyperref[DC18-47]{verse 47} in section 18.
						\end{compactitem}
				\end{compactitem}
			\item DC 19
				\begin{compactitem}
					\item \hyperref[DC19-4]{DC 19:4--12}
						\begin{compactitem}
							\item This is one of the most important passages in the scriptures when it comes to understanding the principles of revelation, because it very clearly explains that a certain phrase in the BOM does \textit{not} mean what we would think it means, based on the most natural way of understanding the words. This is extremely important.
							\item The Lord explains the meaning of the terms in question in 10--12, and that's important. But I think the most important verse of all here is \hyperref[DC19-7]{verse 7}, where the Lord explains that he uses the phrase ``eternal damnation'' (and others) for a pragmatic purpose: with a somewhat harsh or even scary phrase, he wants to impress on people the seriousness of the situation, and work in them a desire to repent. Note what this means: in the revelations, including the BOM, the Lord is choosing words and terms for pragmatic purposes, \textit{even when, on a natural interpretation, those words and phrases communicate a \textbf{false} idea of some aspect of doctrine or the plan of salvation}. This is a principle that is absolutely fundamental to the epistemology of revelation.
							\item See my notes here, but also my notes on those verses too.
							\item See also \hyperref[DC63-53]{DC63-53} (the verse and below in this list).
						\end{compactitem}
					\item \hyperref[DC19-38]{DC 19:38}
				\end{compactitem}
			\item \hyperref[DC20-8]{DC 20:8--12}
						\begin{compactitem}
							\item See my notes on these verses---it's support for the BY quotes about the translation of the BOM.
						\end{compactitem}
			\item \hyperref[DC20-35]{DC 20:35}
			\item \hyperref[DC21-4]{DC 21:4--5}
			\item \hyperref[DC21-9]{DC 21:9}
			\item \hyperref[DC24-5]{DC 24:5--6}
			\item \hyperref[DC25-4]{DC 25:4}
			\item \hyperref[DC25-9]{DC 25:9}
			\item \hyperref[DC27-5]{DC 27:5}
				\begin{compactitem}
					\item The BOM is a revelation, and so the principles which govern the giving and receiving of revelation apply to it.
				\end{compactitem}
			\item \hyperref[DC28]{DC 28}
				\begin{compactitem}
					\item The whole saga of this story with Page and the seer-stone is important for understanding revelation on Earth.
					\item Important verses include 1--4, 5, 7, 8, 11--12.
					\item 8 in particular is very important---you can have revelations, but don't write them as commandments.
					\item[{\hyperref[DC28-7]{DC 28:7}}]
						\begin{compactitem}
							\item This verse belongs in a group that also includes \hyperref[DC35-18]{DC 35:18}, \hyperref[DC64-5]{DC 64:5}, and \hyperref[DC84:19]{DC 84:19}, where the revelations are talking about the ``keys'' and ``mysteries'' together.
						\end{compactitem}
					\item In 11--12, notice that the Lord doesn't say, ``Seer stones are dumb.'' He just says, those things that you did get from there, aren't of me. 
				\end{compactitem}
			\item \hyperref[DC29]{DC 29}
				\begin{compactitem}
					\item See \hyperref[DC29-33]{verses 33--34} especially, where the Lord explains that he is using certain language about beginnings and endings ``that [we] may naturally understand.'' But that language in fact does not really apply to him. Similar to the stuff at the beginning of \hyperref[DC19]{DC 19}. For verse 34 in 29, the distinction between ``spiritual'' and ``temporal'' doens't really mean anything to the Lord. Note also that in \hyperref[DC29-35]{verse 35} the Lord specifically says that the commandments are not ``natural''---but earlier he said that he said something in a certain way so that we would ``naturally understand,'' but since they aren't ``natural,'' then our understanding is going to be wrong if we stay at that level (note that this line of reasoning requires a bit of equivocation with the terms ``natural'' and ``naturally'').
				\end{compactitem}
			\item \hyperref[DC30-2]{DC 30:2--3}
			\item \hyperref[DC32-4]{DC 32:4}
				\begin{compactitem}
					\item There are various things to note in this verse---one is the ``pretend to no other revelation'' part, and a second is the ``pray always so God will unfold the revelations'' part. But there's another part here about things that are written---there are other places in the scriptures besides here that particular emphasis is given to that which is written down. I think there are divine engineering reasons why this is---it's because we have a finite memory and so we need things written down for us, or we need to write them down to remember them. So if something is written down that makes it especially important. I should check this out elsewhere.
				\end{compactitem}
			\item \hyperref[DC33-16]{DC 33:16}
				\begin{compactitem}
					\item The scriptures are given of God ``for our instruction.'' What does this mean? And connect to other uses of ``instruction''?
				\end{compactitem}
			\item \hyperref[DC34-HB]{Historical background to DC 34}, in the JSP, the bit about how the revelation was received
			\item \hyperref[DC35]{DC 35}
				\begin{compactitem}
					\item \hyperref[DC35-17]{17--19}, what JS received
					\item \hyperref[DC35-20]{20}, how the scriptures are ``in mine own bosom,'' for a particular purpose, the ``salvation'' of God's elect; this also relates to translation, because Rigdon is going to do that
					\item \hyperref[DC35-23]{23}, seems important but I'm not sure how yet
				\end{compactitem}
			\item \hyperref[DC36]{DC 36}
				\begin{compactitem}
					\item \hyperref[DC36-2]{Verse 2}, where a person receives the Comforter ``which shall teach...the peaceable things of the kingdom''
				\end{compactitem}
			\item \hyperref[DC38]{DC 38}
				\begin{compactitem}
					\item \hyperref[DC38-2]{Verse 2}, which may relate to \hyperref[DC130-7]{130:7}
					\item \hyperref[DC38-7]{Verses 7--8}
					\item \hyperref[DC38-13]{Verse 13}---though what exactly is the ``mystery'' the Lord is sharing here?
					\item \hyperref[DC38-29]{DC 38:29--30}---it the Lord kind of saying, you need to learn to discern these things on your own? Let them be revealed to you so that you can understand them?
					\item Another issue here, relating to the whole section, is the connection between this section (and possibly others from around this time) to \hyperref[MOSES-7]{Moses 7} (and possibly other parts of that early JST material, including Moses 6 and 8). They are going on around the same time; they influence each other; they build on each other. This is unavoidable.
				\end{compactitem}
			\item \hyperref[DC39]{DC 39}
				\begin{compactitem}
					\item \hyperref[DC39-6]{DC 39:6}; note comparison to \hyperref[2NEPHI-32-5]{2 Nephi 32:5}
					\item \hyperref[DC39-16]{Verse 16}, how God takes our prayers into account in his actions?
				\end{compactitem}
			\item \hyperref[DC42]{DC 42}
				\begin{compactitem}
					\item \hyperref[DC42-3]{DC 42:3} (there are other examples of this principle elsewhere, like maybe \hyperref[DC84-1]{DC 84:1})
					\item \hyperref[DC42-8]{DC 42:8--9}
					\item \hyperref[DC42-13]{DC 42:13--14}
					\item \hyperref[DC42-15]{DC 42:15}
					\item \hyperref[DC42-16]{DC 42:16}
						\begin{compactitem}
							\item ``As seemeth me good''---related to pragmatism
						\end{compactitem}
					\item \hyperref[DC42-56]{DC 42:56--59}
					\item \hyperref[DC42-61]{DC 42:61}
					\item \hyperref[DC42-65]{DC 42:65}
						\begin{compactitem}
							\item Notice the connection to \hyperref[ALMA-12-9]{Alma 12:9--11}.
						\end{compactitem}
					\item \hyperref[DC42-68]{DC 42:68}
					\item An important thing about all these DC 42 verses is what section of the section they are in---remember it was given in parts in response to specific questions. So you would want to look at what each question was that prompted the response, and ask how these epistemology-related points fit in with that question.
				\end{compactitem}
			\item \hyperref[DC43]{DC 43}
				\begin{compactitem}
					\item Like section 28 above, this whole section is important for revelation, including the historical background. It has to do with a woman named Hubble who claimed to be receiving revelations at the time.
					\item I think it's important to note that there's a lot of circularity here---you've got someone receiving a revelation which then says that no one else can receive revelations.
					\item The whole section is important, including the opening verses, even though I don't mention them below.
					\item \hyperref[DC43-5]{Verse 5}, ``revelations'' and ``commandments'' have some special status, above just regular ``teachings''; see something like \hyperref[DC88-118]{DC 88:118}.
					\item \hyperref[DC43-6]{Verse 6}, I think there's some possible tension here with things like \hyperref[MORONI-7-13]{Moroni 7:13}; it could be that, given the actual situations of the saints at this time, it was not good for them to believe Hubble's ``revelations,'' because they were too fragile or just not experienced enough or something.
					\item \hyperref[DC43-7]{Verse 7}, so the ``gate'' metaphor here is being used as a way of talking about doing things in the correct manner or following the correct rules or guidelines for how to do something.
					\item \hyperref[DC43-8]{DC 43:8--10} is a very important set of verses; see the JSP version especially, which and has very important differences, such as the beginning of \hyperref[DC43-9]{verse 9}. These are crucial.
					\item And notice that verse 10 contains the ``if you don't do it, it will be taken away'' principle, very similar to \hyperref[2NEPHI-28-30]{2 Nephi 28:30}.
					\item \hyperref[DC43-11]{DC 43:11--14}, very interesting principles here, about how JS is going to be the one who reveals the ``glories'' and ``mysteries'' of the kingdom, and everyone else has to provide for him to make that possible.
					\item \hyperref[DC43-15]{DC 43:15--16}, more stuff about receiving guidance of the spirit.
					\item \hyperref[DC43-25]{DC 43:25}, the Lord uses many means to speak.
				\end{compactitem}
			\item \hyperref[DC44]{DC 44}
				\begin{compactitem}
					\item \hyperref[DC44-2]{Verse 2}, about gathering and receiving the Spirit (note the attitude or the behavior that those gathering should have)
				\end{compactitem}
			\item \hyperref[DC45]{DC 45}
				\begin{compactitem}
					\item Verse 46 is related to the seeing/knowing thing.
				\end{compactitem}
			\item \hyperref[DC46]{DC 46}
				\begin{compactitem}
					\item \hyperref[DC46-1]{Verses 1--2}, both the ``profit and learning'' part and the meetings part
					\item \hyperref[DC46-7]{Verse 7}, very interesting verse, especially the ``doctrines'' part
					\item \hyperref[DC46-8]{Verses 8--9}, which discuss the purpose of spiritual gifts, or for what they are given
					\item \hyperref[DC46-13]{Verses 13--14}, extremely important for many reasons.
						\begin{compactitem}
							\item How does the Holy Ghost teach them this, in verse 13? By showing them, I think; they \textit{see} it.
							\item Compare this function of the Holy Ghost with \hyperref[DC21-9]{DC 21:9}.
							\item Note the beginning of verse 14---it says, ``to others,'' not ``to all others.'' So it seems like it means ``to some others'' or ``to certain others.''
							\item Notice the very apparent difference between ``know'' in the one verse and ``believe'' in the other.
						\end{compactitem}
					\item \hyperref[DC46-15]{Verse 15}---note that this verse gives you a parallel to revelation. It's talking about God changing what he does in response to different circumstances in order to minister to people, and he does the same thing with revelation. It's parallel to the stuff BY says about translation.
						\begin{compactitem}
							\item In fact, you may even be able to see that idea about revelation as just a special case of this one, which would make this principle even more fundamental.
							\item That is an idea that is worth exploring.
						\end{compactitem}
					\item \hyperref[DC46-22]{Verse 22}---note that some of the spiritual gifts listed here are relevant for epistemology, such as the one in this verse, the gift to prophesy. And notice that it's also relevant that different gifts are given to different people.
					\item \hyperref[DC46-28]{Verses 28--33}---important for asking, also connected to the temple and asking, probably.
				\end{compactitem}
			\item \hyperref[DC47-1]{DC 47:1, 4}
				\begin{compactitem}
					\item Don't forget verse 4---the Comforter is involved here
				\end{compactitem}
			\item \hyperref[DC48-4]{DC 48:4--5}
				\begin{compactitem}
					\item This could be an example of the need to ``study it out in your own mind,'' which includes gathering information---it could be that the added perspectives and knowledge of other people played a role in getting JS and others to the revelation and understanding that they needed. 
					\item Would need to look at the particulars of some of the history for this.
				\end{compactitem} 
			\item \hyperref[DC49]{DC 49}, the whole thing with the Shakers is interesting
				\begin{compactitem}
					\item \hyperref[DC49-2]{Verse 2}, interesting in light of what JS later says about people saying things like, ``That's enough truth for me, thank you.'' Goes with stuff like \hyperref[2NEPHI-28-30]{2 Nephi 28:30}. See also places like \hyperref[1CORINTHIANS-13]{1 Corinthians 13}, where you get this ``in part'' phrase too.
					\item \hyperref[DC49-4]{Verse 4}, more ``reasoning'' stuff; see \hyperref[REASON]{Reason}.
				\end{compactitem}
			\item \hyperref[DC50]{DC 50}
				\begin{compactitem}
					\item See HDDC 647--648 for PPP's discussion of how this revelation was received; this is an important topic in itself, and important more generally, not just for this section.
					\item This section is also important generally, for the ``discerning of spirits'' stuff, which comes up a lot throughout JS's whole career and is important for the endowment ceremony stuff as well.
						\begin{compactitem}
							\item Connect here to \hyperref[DC46-7]{DC 46:7} and maybe other verses around there.
							\item Also, a connection I hadn't really seen before now, connect to \hyperref[2NEPHI-2-29]{2 Nephi 2:29} and stuff in the BOM about spirits, from places like \hyperref[ALMA-34-34]{Alma 34:34}.
						\end{compactitem}
					\item \hyperref[DC50-1]{Verse 1}, an example of agreement among people and the affect it has on receiving revelation
					\item \hyperref[DC50-9]{Verses 9--36}---here begins one of the most remarkable passages in all of the DC. I have a lot of notes below and in the actual section itself; I'm sure I didn't get everything that is worth looking at.
						\begin{compactitem}
							\item[9] Think of truth as goodness here, and righteousness as lawful behavior in the manner of \hyperref[DIKAIOSUNE]{\grl{δικαιοσύνη}}.
							\item[10] The purpose of reasoning together is for us to \textit{understand}.
							\item[10--12] One of the most important sets of verses anywhere for the epistemology of revelation---the Lord will reason with us \textit{as a man}. Compare \hyperref[2NEPHI-31-3]{2 Nephi 31:3} and \hyperref[DC1-24]{DC 1:24}. Need to think more about the implications of these verses.
							\item[14] The Comforter is ``sent forth to teach the truth.'' Compare to truth/goodness stuff and other scriptures on the Holy Ghost/Comforter. Compare verse 17 too.
							\item[15] Need to think about what these ``spirits'' are---see above, on the ``discerning of spirits'' thing at the beginning of the notes on this section.
							\item[17--22] There's something about sharing the ``spirit of truth'' that makes preaching possible. See that phrase \hyperref[SPIRIT-PHRASES]{here}. You have to preach and receive by the ``spirit of truth.'' Verse 21 in particular also makes me think of the ``sower sows the word'' parable given by Christ. Also notice that this way of doing it leads to mutual understanding (verse 22).
							\item[22--24] Extremely important principles---lots of connections here about light and darkness. See also the connection to truth and goodness, with ``edify''---23 is particularly important for this. Also, I have to see the principles in 24 as giving another version of the same prinfciple taught elsewhere, such as \hyperref[DC93-28]{DC 93:28} and \hyperref[2NEPHI-28-30]{2 Nephi 28:30}. It would be useful to round all those principles up into the same place in order to get the full picture.
							\item[25] Knowledge and the truth chase darkness from us.
							\item[27] This verse describes the ``perfection'' or the ``perfect day'' that results from getting so much light. This connects to laws, priesthood, and to the opening of \hyperref[DC88]{DC 88}.
							\item[28--29] Personal righteousness is a big factor here.
							\item[31] There is some stuff about spirits in and around this part that could be important, also possibly related to the endowment ceremony and later developments in 1842.
							\item[34--35] Good verse, especially 35. More results of getting light and keeping it and continuing in God.
						\end{compactitem}
					\item And don't forget about \hyperref[DC50-40]{verse 40}!
				\end{compactitem} 
			\item \hyperref[DC51]{DC 51}
				\begin{compactitem}
					\item As discussed in HDDC 665, Orson Pratt, in late June 1874, gave a talk at a meeting where he spoke about the process of Joseph receiving this revelation---I don't know the original conference minutes source, but see MS 36:497--499. The specific bit about this revelation is in the last two paragraphs on 498, but the whole thing is actually important for the epistemology of revelation.
					\item \hyperref[DC51-16]{Verses 16--17}, where the Lord says he isn't going to say how long the people are going to be in a particular place, they just need to assume it's going to be years and get to work, and it will be good for them. This is also related to gospel pragmatism and truth as goodness. Compare to Nephi getting to the coast and being like, what do we do now? And they stay there for eight years. See also the stuff in 52 below. Compare \hyperref[DC64-21]{DC 64:21}.
				\end{compactitem}
			\item \hyperref[DC52]{DC 52}
				\begin{compactitem}
					\item \hyperref[DC52-2]{Verse 2}---note how this goes along with \hyperref[DC51-16]{DC 51:16--17}, where the Lord is taking that pragmatic approach. This is a good example of line by line, precept by precept, and this would be a good essay topic to follow these revelations during this period as giving one step at a time. Compare to Nephi getting one step at a time. See also \hyperref[DC52-4]{verses 4--5}.
					\item \hyperref[DC52-9]{Verse 9}
					\item \hyperref[DC52-14]{Verses 14--19}---MORE stuff about discerning spirits, here the Lord is giving a ``pattern'' for making decisions about them.
						\begin{compactitem}
							\item There is a lot going on in these verses that needs to be looked at.
							\item Look at 17, for example---is there a difference between ``revelations'' and ``truths''? Why would the Lord mention them in a series like this if they are identical? Very interesting.
							\item Lots of elements to this ``pattern.''
						\end{compactitem}
				\end{compactitem}
			\item \hyperref[DC56]{DC 56}
				\begin{compactitem}
					\item \hyperref[DC56-2]{Verses 2--4}---you can see this as an expansion or commentary on the idea of receiving less light as you disobey. Note the timeframe in verse 3 is up to the Lord, and in verse 4 that's clear too, with ``seemeth [him] good.'' 
					\item In this whole section really, notice how the Lord is revoking commandments in response to different things. Why didn't he just tell them what to do using his all-knowing powers? Why didn't he just tell them the final thing that they would end up doing? This is relevant for many sections in this middle part of the DC.
					\item This goes with other sections around here, including 51 and 52 above.
				\end{compactitem}
			\item \hyperref[DC57]{DC 57}
				\begin{compactitem}
					\item Interesting to see all the times that the word ``wisdom'' is used here. The Lord gives very specific and very direct instruction here, as opposed to more general guidelines elsewhere, and he calls it ``wisdom'' several times.
					\item Also verse 13 is interesting.
				\end{compactitem}
			\item \hyperref[DC58]{DC 58}
				\begin{compactitem}
					\item \hyperref[DC58-3]{Verses 3--4}, especially the ``natural eyes'' part. 
					\item \hyperref[DC58-5]{Verse 5}, remember in order to receive more
				\end{compactitem}
			\item \hyperref[DC59]{DC 59}
				\begin{compactitem}
					\item \hyperref[DC59-3]{Verses 3--4}
				\end{compactitem}
			\item \hyperref[DC61]{DC 61}
				\begin{compactitem}
					\item \hyperref[DC61-22]{Verse 22}, pragmatic
					\item \hyperref[DC61-27]{Verses 27--28}---verse 27 could be pointing to the relationship in the priesthood between power to do things and power to know things (see the 6 April 1837 from JS at \hyperref[PRIESTHOOD]{Priesthood}). 
				\end{compactitem}
			\item \hyperref[DC62]{DC 62}
				\begin{compactitem}
					\item \hyperref[DC62-5]{Verse 5}, pragmatic (see also verses 7--8)
					\item \hyperref[DC62-6]{Verse 6}, the Lord can't lie about the promises he makes
				\end{compactitem}
			\item \hyperref[DC63]{DC 63}
				\begin{compactitem}
					\item \hyperref[DC63-7]{Verse 7--12}---there's something going on here with the experience of seeing a sign but it not leading to the right outcome, or change of behavior, or something (see also verse 11). See verses 8--12 too, especially the part about how ``faith cometh not by signs.'' This would be important to work out---what does it mean, and why is it true? And what does it mean that ``signs follow those that believe'' (9)?
					\item Pragmatic bit with \hyperref[DC63-12]{verse 12}.
					\item \hyperref[DC63-16]{Verse 16}, pretty clearly linking righteousness with having the Spirit (or at least getting it taken away or something).
					\item \hyperref[DC63-23]{Verse 23}, absolutely fundamental. One of the best examples of this principle in all the scriptures. Goes with \hyperref[ALMA-12-10]{Alma 12:10} too.
						\begin{compactitem}
							\item More and more I am sensing a relationship between the priesthood, righteousness (as lawful behavior), and revelation---they all go together. They are all different aspects of the same thing, or at least they are very intimately related in some ways.
							\item Also related to the \hyperref[NEWBIRTH]{new birth}.
							\item On this point, note that many discussions of the priesthood include epistemological points about receiving revelation or having the right to it.
						\end{compactitem}
					\item \hyperref[DC63-26]{Verse 26}, definitely shows the pragmatic side, this is very important. (Note \hyperref[DC63-40]{verse 40} too.)
					\item \hyperref[DC63-41]{Verse 41}, a specific power of discernment mentioned here, for a very specific purpose.
					\item \hyperref[DC63-53]{Verse 53}---now this is important. It's another example of what goes on in the early part of \hyperref[DC19]{DC 19}. The Lord is explaining something but he explicitly says that he is ``speaking after the manner of the Lord''---so we would be wrong to assume that we can bring our normal understanding of concepts or terms to what he is saying. He's speaking in a different language, or with a different background of assumptions; incredibly important. See also 64:24.
					\item \hyperref[DC63-64]{Verses 64--65}, not totally sure what to make of this, at least in verse 64.
				\end{compactitem}
			\item \hyperref[DC64]{DC 64}
				\begin{compactitem}
					\item \hyperref[DC64-5]{DC 64:5}, also important for prophets
					\item \hyperref[DC64-16]{DC 64:16}
					\item \hyperref[DC64-20]{DC 64:20}, pragmatic (see also verses 27--28)
					\item \hyperref[DC64-21]{DC 64:21}, note the specific timeframe given here, compared to \hyperref[DC51-17]{DC 51:17}
					\item \hyperref[DC64-23]{Verses 23--25}, a couple things to note:
						\begin{compactitem}
							\item The Lord tells us something about time: ``now it is called today.'' So we're using concepts in a figurative way here. Note the use of ``tomorrow'' in verse 24, in comparison to ``today.'
							\item Here's another example of ``speaking after the manner of the Lord,'' in verse 24. Extremely important. See 63:53 too.
						\end{compactitem}
					\item \hyperref[DC64-34]{Verse 34}, getting the blessings of Zion (part of which will be the Spirit and revelation)
				\end{compactitem}
			\item \hyperref[DC65]{DC 65}
				\begin{compactitem}
					\item \hyperref[DC65-2]{Verse 2}, giving the keys involves both the priesthood and revelation; think of \hyperref[MATTHEW-16-18]{Matthew 16:18} too.
				\end{compactitem}
			\item \hyperref[DC66]{DC 66}
				\begin{compactitem}
					\item \hyperref[DC66-2]{Verse 2}, the everlasting covenant and the fulness of the gospel; glories revealed in the last days. Why in the last days? I would say, at least partially due to ventricular theology reasons.
					\item \hyperref[DC66-7]{Verse 7}, another use of ``reasoning''
				\end{compactitem}
			\item \hyperref[DC67]{DC 67}
				\begin{compactitem}
					\item \hyperref[DC67-1]{DC 67:1--9}---In connection with this, see the HDDC discussion of DC 20 in the \hyperref[DC20-HB]{historical background} section, especially HDDC 292. See also the \hyperref[DC67-HB]{historical background} for section 67 and for \hyperref[DC1-HB]{section 1}, which are extremely important, including the ``testimony'' of the elders about the revelations (discussed at 67). Note the discussion in the historical background to 67 about the conversation on ``Revelations and language.'' HDDC 21--26 has more.
						\begin{compactitem}
							\item Verse 4, the Lord himself is giving a testimony of their truth.
							\item There's really a lot going on here, especially starting in verse 5. That's where it really gets interesting.
							\item Notice that the language being spoken of in verse 5 is \textit{JS}'s language. This is because the revelations had to be given in JS's language, exactly as the Lord stated they would, in \hyperref[2NEPHI-31-3]{2 Nephi 31:3}. And for this exact reason, the Lord says in the preface to the DC, \hyperref[DC1]{section 1}, that he speaks in the language of the people (verse 24). (If the JSP dating is correct, it's possible that DC 1 came before 67.)
							\item 6--9 is the test the Lord is proposing; it's an interesting test. Verse 7 delivers the promise of the test.
						\end{compactitem}
					\item \hyperref[DC67-1]{DC 67:1, 3}---Notice that the Lord says they ``endeavored to \textit{believe},'' rather than \textit{do} something; they tried to hold on to that hope and faith.
					\item \hyperref[DC67-10]{10--14}---extremely important verses, especially 10. One of the most important verses in the entire DC probably, in terms of epistemology stuff. We will see Christ! See notes at verse 10 for references to the phrases involving ``mind.'' Very important to think about this kind of perception and what it means.
						\begin{compactitem}
							\item These four verses are like a statement of everything you want from righteousness. Just beautiful and powerful.
							\item Note the ``in mine own due time,'' at the end in verse 14. We don't get to choose this.
						\end{compactitem}
					\item Verses 10--13 are important, especially 10--11; I think these tell us what it means to perceive or experience with the ``spiritual mind.'' It means to be quickened, or to have the Spirit of God acting on us in order to influence how and what we perceive.
					\item Verse 14 is useful too; the whole section is very useful for epistemology of revelation stuff.
				\end{compactitem}
			\item \hyperref[DC68]{DC 68}
				\begin{compactitem}
					\item Very important in this section is the addition made of 15--21 to the later EMS-R version of the section; that wasn't in the original. So this is a case study of going back and making major changes, and what does that say about these revelations? What does it say about an evolving understanding? Very important to work this out. Use \hyperref[DC107]{section 107} to think about this, because the stuff added to 68 comes after that. See the \hyperref[DC68-HB]{historical background} to 68.
					\item \hyperref[DC68-3]{Verses 3--4}---not sure what to do with these ones, but verse 4 in particular is really important. Could relate also to pragmatic stuff. Notice that it doesn't say that any of this stuff is ``true.''
					\item \hyperref[DC68-11]{Verse 11}, knowing signs
					\item \hyperref[DC68-21]{Verse 21}, following after many important verses about the priesthood (which were not in that form in the original), but here it says you could ascertain your lineage ``by revelation.'' Seems to be referring to patriarchal blessings, but it also talks about being a ``literal descendent'' (verse 18), and if that's what we're talking about here, then I don't see how different people in different families could have different lineages. If that's possible then we aren't talking about ``literal descendents.''
					\item \hyperref[DC68-25]{Verse 25}, important, note the ``to understand'' part. Also important for accountability.
				\end{compactitem}
			\item \hyperref[DC70]{DC 70}
				\begin{compactitem}
					\item \hyperref[DC70-1]{Verses 1--2}---note here that the Lord explicitly says that he is giving a commandment.
					\item \hyperref[DC70-12]{Verses 12--15}, mostly 12--14; the ``abundance'' these people get is in the ``manifestations of the Spirit.'' Sounds pretty amazing. Note also 14, about the relationship between these manifestations and your attitude toward your temporal possessions. 
				\end{compactitem}
			\item \hyperref[DC71]{DC 71}
				\begin{compactitem}
					\item \hyperref[DC71-1]{Verse 1}, you get a \textit{portion} of the ``spirit and power'' in order to expound the mysteries
					\item \hyperref[DC71-4]{Verse 4}, prepare the way in the minds of people for future commandments, because whatever we tell them later is going to be built on top of what we tell them now. 
					\item \hyperref[DC71-5]{Verses 5--6}, a very interesting version of the ``obey/get more'' principle, relating to power; connect back to verse 1 here too. If you receive wisdom, you will get \textit{power}. Connects to the ``knowledge is power'' principle too, in JS; see \hyperref[KNOWLEDGE]{Knowledge}.
					\item[DC71-8]{Verse 8}, reasoning and reasons.
				\end{compactitem}
			\item \hyperref[DC73]{DC 73}
				\begin{compactitem}
					\item \hyperref[DC73-5]{Verse 5}, one step at a time
				\end{compactitem} 
			\item \hyperref[DC74]{DC 74}
				\begin{compactitem}
					\item Interesting story behind this section---it's an explanation of a verse in the NT, \hyperref[1CORINTHIANS-7-14]{1 Corinthians 7:14}. You could look at how JS is explicating it and what that means for the revelation he receives.
					\item \hyperref[DC74-5]{Verse 5}, crucial verse on the difference between a commandment ``of the Lord'' and one of a person (even an apostle). How can we tell the difference?
				\end{compactitem}
			\item \hyperref[DC75]{DC 75}
				\begin{compactitem}
					\item \hyperref[DC75-1]{Verse 1}---the Lord speaks by ``the voice of his spirit.'' Could be an interesting phrase to connect elsewhere.
					\item Verse 4, proclaiming the truth \textit{according to} the revelations and commandments.
					\item \hyperref[DC75-10]{Verse 10}, we get the Comforter to teach us ``all things that are expedient''; has to do with pragmatism and also context
					\item \hyperref[DC75-27]{Verse 27}, similar to verse 10
				\end{compactitem}
			\item \hyperref[DC76]{DC 76}
				\begin{compactitem}
					\item The story of how this section was received is very important for the epist stuff too---see the accounts in the \hyperref[DC76-HB]{historical background}. Lots to talk about there. Note verse 15---it's while they were doing the translation.
					\item \hyperref[DC76-2]{Verse 2}, we can't find out the extent of the Lord's doings.
					\item \hyperref[DC76-5]{Verses 5}, righteousness is the key
					\item \hyperref[DC76-7]{DC 76:7--10}, extremely important verses, but make sure you pick up on verses 5--6 too, since those describe which people we are talking about here, or which people are going to receive this reward
					\item \hyperref[DC76-12]{DC 76:12}, it's ``by the power of the spirit'' that their eyes were opened (see also \hyperref[DC76-19]{verse 19})
					\item \hyperref[DC76-45]{DC 76:45--48}, about the torment and who gets to know about it
					\item \hyperref[DC76-89]{89}, more about who can know this stuff
					\item \hyperref[DC76-114]{DC 76:114--118}---super important. Full of good stuff.
				\end{compactitem}
			\item \hyperref[DC77]{DC 77}
				\begin{compactitem}
					\item This whole explication of the NT is relevant for thinking about the translation stuff and for revelation.
					\item See verse 2 in particular, with the relation between spiritual and temporal things; this is only a symbolic relationship, remember, used to talk about the imagery in the scriptures.
				\end{compactitem}
			\item \hyperref[DC78]{DC 78}
				\begin{compactitem}
					\item \hyperref[DC78-10]{DC 78:10}, how Satan gets us away from what is true
					\item \hyperref[DC78-17]{Verses 17--18}, we're children in our understanding but will be led along
					\item Note the ``Ahman'' stuff in this section, which isn't in the original; this is a good example of later revisions.
				\end{compactitem}
			\item \hyperref[DC79]{DC 79}
				\begin{compactitem}
					\item \hyperref[DC79-2]{verse 2}, the Comforter does the teaching
				\end{compactitem}
			\item \hyperref[DC80]{DC 80}
				\begin{compactitem}
					\item \hyperref[DC80-3]{verse 3}, pragmatism, sometimes the place doesn't matter. 
					\item \hyperref[DC80-4]{verse 4}, distinction between hearing, believing, and knowing to be true
				\end{compactitem}
			\item \hyperref[DC81]{DC 81}
				\begin{compactitem}
					\item The historical introduction has some important material about how this revelation was later seen, given Jesse Gause's later actions.
					\item See \hyperref[DC81-4]{verse 4}, which talks about doing the ``greatest good.'' This is a way of seeing the purpose of the instructions that the Lord gives us.
				\end{compactitem} 
			\item \hyperref[DC82]{DC 82}
				\begin{compactitem}
					\item \hyperref[DC82-3]{verses 3--4}, very important, a general statement of this principle. Reminds me of \hyperref[DC93]{DC 93} and the verses about agency. Verse 4 is also really important.
					\item Verse 7? Not sure if it bears on epistemology.
					\item \hyperref[DC82-8]{verses 8--10}, all three are very important in different ways. Things ``turning to us'' for our salvation is about as clear a description of the kind of ``liquid luck'' idea that you're going to get.
				\end{compactitem}
			\item \hyperref[DC84]{DC 84}
				\begin{compactitem}
					\item \hyperref[DC84-1]{verses 1--2}, the story behind this revelation, and how it's given to a group of people ``as they united their hearts and lifted their voices on high.'' It's also interesting that it's explicitly called a ``revelation'' in verse 1.
					\item \hyperref[DC84-19]{verses 19--22}, about the relationship between the priesthood and knowledge.
					\item \hyperref[DC84-43]{verses 43--54}, lots of good material here about living by the Spirit but also other things---this block of twelve verses has a lot of different things going on
					\item \hyperref[DC84-57]{verse 57}, on translation and whether the Lord is saying he wrote the Book of Mormon
					\item \hyperref[DC84-60]{verse 60}, very important relationship between ``words'' and ``voice''; note how this changes how we should read earlier verses, like 52, in this section.
					\item \hyperref[DC84-75]{verse 75}, timing of effect
					\item In roughly the second half of this revelation, really starting in the 50's of the verses, Christ begins giving instructions very similar to the ones that he gave to his followers during his time around the Jerusalem area.
					\item \hyperref[DC84-85]{verse 85}, good missionary scripture but also important for epist stuff
					\item \hyperref[DC84-98]{verse 98}, we'll ``see eye to eye''
					\item \hyperref[DC84-106]{verses 106--110}, I'm interested here in the teamwork aspect of this, and how we need everyone to keep the system ``perfect''
					\item \hyperref[DC84-117]{verse 117}, see note there
					\item \hyperref[DC84-119]{verse 119}, see/know, compare \hyperref[DC93-1]{DC 93:1}; think of the knowledge growing over time as we progress
				\end{compactitem}
			\item \hyperref[DC85]{DC 85}
				\begin{compactitem}
					\item \hyperref[DC85-6]{verse 6}, on the ``still small voice.'' See notes there.
					\item \hyperref[DC85-7]{verse 7}, this verse appears to be a direct quotation of something given to JS through the Spirit.
				\end{compactitem}
			\item \hyperref[DC86]{DC 86}
				\begin{compactitem}
					\item I'm not sure what it is, but it seems like there is something significant about the wheat and the tares growing and existing together at once. 
				\end{compactitem} 
			\item \hyperref[DC87]{DC 87}
				\begin{compactitem}
					\item Interesting to see early events in South Carolina, and how those would have affected JS's views on the United States---even back in 1832, as noted in the historical background to section 87, there were good reasons to think that South Carolina was going to be the root of the problem. But the war was still 28 years away, and so that can't take away from this section, which really is marvelous. 
				\end{compactitem}
			\item \hyperref[DC88]{DC 88}
				\begin{compactitem}
					\item \hyperref[DC88-32]{Verse 32}, this verse could apply to epistemology as well, in ``enjoying that which you are willing to receive.''i
					\item \hyperref[DC88-63]{Verse 63}, about Christ drawing near to us
						\begin{compactitem}
							\item Verses \hyperref[DC88-64]{64--68} are also important in connection with this.
						\end{compactitem}
					\item The parable about the laborers, including \hyperref[DC88-46]{46--61}---see my note at \hyperref[DC88-51]{verse 51}.
					\item There are other verses that could work for epistemology too, since some of the section is about receiving light---but it's not always referring to epistemology or knowledge, I think. There's more going on here.
					\item \hyperref[DC88-77]{verses 77--79}					
					\item \hyperref[DC88-118]{verses 118--126}
				\end{compactitem}
			\item \hyperref[DC89]{DC 89}
				\begin{compactitem}
					\item \hyperref[DC89-3]{Verse 3}
					\item \hyperref[DC89-18]{Verses 18--21}
				\end{compactitem}
			\item \hyperref[DC90]{DC 90}
				\begin{compactitem}
					\item \hyperref[DC90-2]{Verses 2--5}, especially the ``oracles'' partsd
					\item \hyperref[DC90-9]{Verse 9}
					\item \hyperref[DC90-11]{Verse 11}, very important verse; connect back to \hyperref[2NEPHI-31-3]{2 Nephi 31:3}
					\item \hyperref[DC90-13]{Verses 13--15}
					\item \hyperref[DC90-24]{Verse 24}
				\end{compactitem}
			\item[{\hyperref[DC91]{DC 91}}]
				\begin{compactitem}
					\item Both the stuff about the Apocrypha at the beginning, and also the instructions for reading the scriptures, which start at verse 4.
					\item The Spirit manifesting ``truth'' and the benefits you get from that are important.
				\end{compactitem}
			\item[{\hyperref[DC92]{DC 92}}]
				\begin{compactitem}
					\item \hyperref[DC92-1]{Verse 1}, ``What I say unto one, I say unto all,'' how often does this apply in the scriptures? How often does it not apply?
				\end{compactitem}
			item[{\hyperref[DC93]{DC 93}}] This whole section is amazing, of course, one of the most important of all for the epistemology of the gospel and revelation; it might be the most important, in fact. The whole thing should be studied carefully but I've tried to break it down into its various parts below for a little guide.
				\begin{compactitem}
					\item \hyperref[DC93-1]{verses 1--3} A statement of a promise---do the five things listed in verse 1, and then you will have a certain experience and come to know certain things. The things you'll come to know are listed in verses 1--3.
					\item \hyperref[DC93-4]{verses 4--5}, which are an explanatory comment either to the mentions of the Son and the Father, in verse 3, or perhaps an unspoken question about the sense in which Christ can be considered both the Father and the Son. It's strange because the term ``Son'' hasn't appeared yet until it's involved in an explanation in verse 4. It could also be that verse 4 in particular is explainin the in-dwelling relationship of the Father and the Son in verse 3. Compare \hyperref[MOSIAH-15-2]{Mosiah 15:2--5}.
					\item \hyperref[DC93-6]{verse 6}, new section, a short introduction to an account of John's testimony regarding Christ. It could be that the account starting in verse 7 is the fulness of that account, as mentioned in verse 6.
					\item \hyperref[DC93-7]{verse 7}, begin John's actual account
					\item \hyperref[DC93-11]{verse 11}, continuing John's account, sort of a new emphasis here or a slightly new direction
					\item \hyperref[DC93-12]{verses 12--14}, a section which explains Christ's progression and why he is called the son of God
					\item \hyperref[DC93-15]{verses 15}, one verse which appears to be about Christ's baptism but it doesn't explicitly say that so it doesn't have to be
					\item \hyperref[DC93-16]{verses 16--17}, a statement of how Christ did end up receiving the fulness of the Father
					\item \hyperref[DC93-18]{verses 18}, a promise that we'll get a fulness of John's record 
					\item \hyperref[DC93-19]{verse 19}, which tells you the purpose of everything that has come before it
					\item \hyperref[DC93-20]{verse 20}, a statement which says that we will progress in the same way Christ does
					\item \hyperref[DC93-21]{verses 21--23}, discussion of Christ's identity along with our own (identity and history)
					\item \hyperref[DC93-24]{verses 24--25}, an important statement about truth; pairs with verse 25
					\item \hyperref[DC93-26]{verses 26--28}, a statement again about our progression
					\item \hyperref[DC93-29]{verses 29--30}, statements in intelligence and truth
					\item \hyperref[DC93-31]{verses 31--32}, statements about agency
				\end{compactitem}
			\item[{\hyperref[DC94-17]{DC 94:17}}] This verse, along with some others in this section, is a good example of ``unfolding'' revelation across time. 
			\item[{\hyperref[DC95]{DC 95}}]
				\begin{compactitem}
					\item[{\hyperref[DC95-1]{DC 95:1--6}}] Especially verse 6.
					\item[{\hyperref[DC95-7]{DC 95:7}}] This is a really weird thing about ``Sabaoth.''
					\item[{\hyperref[DC95-8]{DC 95:8--9}}] Getting endowed with power from on high; it's a promise from the Father
					\item[{\hyperref[DC95-12]{DC 95:12}}] Very interesting verse about the ``love'' of the Father not continuing with us if we don't obey the commandments. See notes at \hyperref[2NEPHI-4-21]{2 Nephi 4:21}.
					\item[{\hyperref[DC95-17]{DC 95:17}}] See notes here; very strange stuff going on.
				\end{compactitem}
			\item[{\hyperref[DC96]{DC 96}}] There are various uses of ``expedient'' in here, as the Lord gives wisdom to ``know how to act'' (1).
			\item[{\hyperref[DC97]{DC 97}}]
				\begin{compactitem}
					\item \hyperref[DC97-1]{verse 1}, the Lord's voice is the voice of his spirit; this could be important for unraveling certain passages elsewhere, including passages about ``living by every word that proceeds from God's mouth''
					\item \hyperref[DC97-1]{verses 1--2}, a good example of asking and then receiving
					\item \hyperref[DC97-5]{verse 5}, PPP is blessed to expound the scriptures and mysteries
					\item \hyperref[DC97-12]{verses 12--14}, building a house in order to instruct people
				\end{compactitem}
			\item[{\hyperref[DC98]{DC 98}}]
				\begin{compactitem}
					\item \hyperref[DC98-11]{verse 11}, live by every word that proceeds forth from the mouth of God; also, ``cleave unto all good''
					\item \hyperref[DC98-12]{verse 12}, not separate from the previous verse, but it does have an elaboration, the idea of ``line upon line.'' Notice the ``herewith'' part.
					\item \hyperref[DC98-20]{verse 20}, we are to observe the ``words of wisdom and eternal life''; pair with 21--22 to see the benefits
				\end{compactitem}
			\item[{\hyperref[DC99-2]{DC 99:2--4}}] About receiving an individual and therefore receiving Christ/the kingdom
			\item[{\hyperref[DC101]{DC 101}}]
				\begin{compactitem}
					\item \hyperref[DC101-22]{verses 22--25}, prepare for a revelation; note the connection to the ancient tabernacle
					\item \hyperref[DC101-25]{verses 25--34}---these verses are describing some sort of millennial period, it seems, and a number of them refer to special revelation or knowledge that will be available during this time. There will be other important things happening during this time, like Satan not being able to tempt people (verse 28).
					\item \hyperref[DC101-95]{verse 95}, the Lord's work described as a ``strange act'' (need to look for biblical parallels)
				\end{compactitem}
			\item[{\hyperref[DC102]{DC 102}}]
				\begin{compactitem}
					\item \hyperref[DC102-1]{verses 1--2}, this high council was specifically appointed by revelation---very practical, emphasis on what needs to get done right now. See also verse 9 in this section.
					\item \hyperref[DC102-23]{verse 23}, getting instruction to clarify matters of doctrine and principle. Note that getting revelation on these sorts of cases seems to be the exception, rather than the rule.
				\end{compactitem}
			\item[{\hyperref[DC103]{DC 103}}]
				\begin{compactitem}
					\item \hyperref[DC103-4]{verses 4--8}, goes toward the ``obey and receive more'' idea
				\end{compactitem}
			\item[{\hyperref[DC104]{DC 104}}]
				\begin{compactitem}
					\item \hyperref[DC104-80]{DC 104:80--81}, extremely interesting description of how the Spirit affects other people---it has some kind of affect such that it is able to change the thoughts of other people in certain circumstances
				\end{compactitem}
			\item[{\hyperref[DC105]{DC 105}}]
				\begin{compactitem}
					\item \hyperref[DC105-40]{DC 105:40}, ``according to the voice of the Spirit which is in you''
				\end{compactitem}
			\item[{\hyperref[DC107]{DC 107}}]
				\begin{compactitem}
					\item \hyperref[DC107-18]{verses 18--20}, relationship of the priesthood to revelation
					\item \hyperref[DC107-30]{verses 30--31}, making decisions in groups with a certain attitude leads to knowledge
					\item \hyperref[DC107-71]{verse 71}, about having knowledge by the spirit of truth
				\end{compactitem}
			\item[{\hyperref[DC109]{DC 109}}]
				\begin{compactitem}
					\item An interesting thing here with epistemology is that this ``revelation'' appears to have been written by multiple people---see Oliver Cowdery's journal entry in the history background. Speaks to the constructive nature of revelation---so the ``construction'' here is the mental/cognitive construction, but also a collaborative construction process with others in some cases
					\item \hyperref[DC109-12]{verses 12--13}, about feeling the power and presence of the Lord
					\item \hyperref[DC109-14]{verses 14--15}, about maturing in the gospel and the light you receive as you go through that process. See my notes there and connections to other verses.
				\end{compactitem}
			\item[{\hyperref[DC110]{DC 110}}] The entire section is relevant, because of the nature of the manifestation, but a few verses are relevant in particular:
				\begin{compactitem}
					\item \hyperref[DC110-1]{verse 1}, about the ``veil'' and ``the eyes of our understanding''
					\item \hyperref[DC110-2]{verse 2}, it's as though a camera is on the Lord and it's showing a view of him from wherever he is, some other place, with a path of amber-colored gold.
					\item \hyperref[DC110-7]{verses 7--8}, the Lord will manifest himself to us and speak to us with his own voice---if we keep the commandments and don't pollute the house
				\end{compactitem}
			\item[{\hyperref[DC111]{DC 111}}] This whole section could be an example of how the Lord uses what we already believe to lead us along to accomplish other purposes
				\begin{compactitem}
					\item \hyperref[DC111-8]{verse 8}, nice description of how the Spirit indicates things to us
					\item \hyperref[DC111-11]{verse 11}, also very important for the pragmatic aspect of revelation and how it's for action, not necessarily for gaining knowledge; the Lord is making use of follies to push us in the right direction eventually
				\end{compactitem}
			\item[{\hyperref[DC112]{DC 112}}]
				\begin{compactitem}
					\item \hyperref[DC112-10]{verse 10}, the Lord God will lead him by the hand---like one step along the way, and the prayers will be answered. Good for thinking about pragmatism. See also verse 19 here, and the idea of an ``effectual'' door.
				\end{compactitem}
			\item[{\hyperref[DC113]{DC 113}}]
				\begin{compactitem}
					\item There's some interesting stuff in the historical background here about how this came about, and how some of it has ``thus saith the Lord,'' and other parts don't.
					\item Also some good stuff here about how JS is using the OT---this is important to look at, I think. Important to see how he's interpreting it and using it as a prophet who is intervening in new and non-traditional ways.
				\end{compactitem}
			\item[{\hyperref[DC121]{DC 121}}]
				\begin{compactitem}
					\item \hyperref[DC121-26]{Verses 26--32}, about revelation and God giving us the Holy Ghost to know things that have never been revealed---see the notes to verse 26 in particular, because the original source has more here that is relevant that is not in the DC section.
					\item \hyperref[DC121-33]{Verse 33}, note that this is its own chunk, need to check out the context in the original
				\end{compactitem}
			\item[{\hyperref[DC123]{DC 123}}]
				\begin{compactitem}
					\item \hyperref[DC123-11]{Verses 11--14}
				\end{compactitem}
			\item[{\hyperref[DC124]{DC 124}}]
				\begin{compactitem}
					\item \hyperref[DC124-1]{Verse 1}, the purpose of raising JS up as a prophet was to show wisdom by means of the weak things of the earth---could be a more general principle about prophets
					\item \hyperref[DC124-4]{Verses 4--5}, special knowledge given about rulers and what is to come, and the emphasis is on the holy ghost
					\item \hyperref[DC124-9]{Verse 9}, God visiting people to change what they believe or how they receive something
					\item \hyperref[DC124-30]{Verses 30--32}---context. In this case, poverty allows the saints to do something that they normally wouldn't be allowed to do.
					\item \hyperref[DC124-40]{Verses 40--46}, the Lord revealing things that haven't been revealed before, here especially in connection with the temple.
					\item \hyperref[DC124-49]{Verse 49}, not really about epistemology, but important for contextual stuff
					\item \hyperref[DC124-91]{Verses 91--97}, important for temple stuff too
					\item \hyperref[DC124-119]{Verses 119--120}, about believing in the BOM and restricting certain activities to those who believe
					\item \hyperref[DC124-123]{Verses 123--126}, about roles in the priesthood and their relation to epistemology
				\end{compactitem}
			\item[{\hyperref[DC128]{DC 128}}]
				\begin{compactitem}
					\item \item\hyperref[DC128-20]{DC 128:20--21}, especially verse 21
				\end{compactitem}
			\item[{\hyperref[DC129]{DC 129}}]
				\begin{compactitem}
					\item There's some extremely specific epistemology stuff here, on discerning spirits; I'm not sure if there are any larger principles, though it is interesting how JS interprets the scriptures that are relevant to this passage.
				\end{compactitem}
			\item[{\hyperref[DC130]{DC 130}}]
				\begin{compactitem}
					\item Verse 3, the personal appearance of God and Christ
					\item Verses 6--11, which are not all one block, but talk about the Urim and Thummim and learning things that are manifest in different ways. These passages would require some study and thought.
					\item \hyperref[DC130-18]{Verses 18--19}, very important verses---see notes, and note that we gain knowledge and intelligence through ``diligence and obedience.'' These verses deserve some more serious thought.
				\end{compactitem}
			\item[{\hyperref[DC131]{DC 131:5--6}}] Interesting verses, and interesting to think of what the ``ignorance'' spoken of is.
			\item[{\hyperref[DC132]{DC 132}}]
				\begin{compactitem}
					\item There's a huge amount to unpack here, including using the principles of the epistemology of revelation to unpack how this revelation could have been received in the first place, since the doctrine itself makes very little sense in the context of some other aspects of the theology.
					\item Verse 3, preparing your heart/receiving a law
					\item Verse 45
				\end{compactitem}
			\item[{\hyperref[DC136]{DC 136}}]
				\begin{compactitem}
					\item \hyperref[DC136-32]{DC 136:32--33}
					\item \hyperref[DC136-37]{DC 136:37}
				\end{compactitem}
			\item[{\hyperref[DC137]{DC 137}}]
				\begin{compactitem}
					\item This entire vision is important, as are all the manifestations had that day by many people.
					\item Note verse 1, the ``body'' thing, referring to the beginning of \hyperref[2CORINTHIANS-12]{2 Corinthians 12}.
				\end{compactitem}
			\item[{\hyperref[DC138]{DC 138}}]
				\begin{compactitem}
					\item Verses 1--5, 11 describe the setting and what was happening when he had this vision; this is a good template to draw some conclusions from.
					\item Verses 28--29, similar to the previous ones
				\end{compactitem}
		\end{compactdesc}
		
	% % -----
	% \subsubsection{Pearl of Great Price} % *PGP
	% % -----
		% \label{EPIST-REVELATION-PGP}
	
		% \begin{tabular}{ll}
			% Total occurrences in the PGP & 
		% \end{tabular}
		
		% \begin{compactdesc}
			% \item[]
		% \end{compactdesc}
		
% % ----------------------
% \subsection{Key Phrases} % *KEYPHRASES *PHRASES
% % ----------------------
	% \label{EPIST-REVELATION-PHRASES}

	% \begin{compactdesc}
		% \item[] \textbf{#times} | list
	% \end{compactdesc}
	
% % ----------------------
% \subsection{Bible Dictionaries} % *DICTIONARIES
% % ----------------------
	% \label{EPIST-REVELATION-BD}

% ----------------------
\subsection{Restoration Usage} % *RESTORATION
% ----------------------
	\label{EPIST-REVELATION-RESTORATION}

	% -----
	\subsubsection{Joseph Smith} % *JS
	% -----
		\label{EPIST-REVELATION-JS}
	
		\begin{compactdesc}
		
			\item[{\hyperlink{EMS-1830-JS-CIR-EARLY-1830-c}{Circa early 1830}}] \& it Pleaseth me that Oliver Cowderey Joseph Knight Hyram Page \& Josiah Stowel shall do my work in this thing yea even in securing the <Copy> right \& they shall do it with an eye single to my Glory that it may be the means of bringing souls unto \sout{me} Salvation through mine only Begotten Behold I am God I have spoken it \& it is expedient in me Wherefor I say unto you that ye shall go to Kingston seeking me continually through mine only Begotten \& if ye do this ye shall have my spirit to go with you \& ye shall have an addition of all things which is expedient in me.
		
			\item[{\hyperlink{EMS-1834-JS-21-APR-1834-a}{21 April 1834}}] Bro. Joseph Smith Jun\supers{r.} read the 2\supers{nd.} chapter of the prophecy of Joel \& took the lead in prayer; after which, he commenced addressing the congregation, as follows. It is very difficult for us to communicate to the churches all that God has revealed to us, in consequence of tradition; for we are differently situated from any other people that ever existed upon this earth: Consequently those former revelations cannot be suited to our condition, because they were given to other people who were before us;
		
			\item[{\hyperref[EMS-1839-JS-CIR-26-JUN-CIR-4-AUG-1839-A-WR]{Between circa 26 June and circa 4 August 1839--A}}] The devil may appear as an angel of light. Ask God to reveal it. if it be of the Devil. he \sout{it be} will \sout{of} flee from you. if of God he will manifest himself or make it manifest. we may come to Jesus <\& ask him> he will know all about it.--- If he comes to a little child, he will adapt himself to the <Language \&> capacity of a little child
			
			\item[{\hyperlink{EMS-1842-JS-28-APR-1842-g}{28 April 1842}}] He said, if God has appointed him, and chosen him as an instrument to lead the church, why not let him lead it through? Why stand in the way, when he is appointed to do a thing? Who knows the mind of God? Does he not reveal things differently from what we expect?
			
			\item[{\hyperlink{EMS-1843-JS-9-JUL-1843-WR-b}{9 July 1843 WR-b}}] one [of] the grand fundamental principles of Mormonism <is> to receivee thruth [truth] let it come from where it may.
			
		\end{compactdesc}
	
	% -----
	\subsubsection{Brigham Young} % *BY
	% -----
		\label{EPIST-REVELATIO 	N-BY}
	
		\begin{compactdesc}
		
			\item[18 April 1844, JSP-AR:119--121] We can form to ourselves independant of the word of the Lord the best system of government on the earth; but after all this, when we have done all the Lord will make it just right. He can form a constitution by which he is willing to be governed. He is willing to be ruled by the means which God will appoint. He dont believe we can adopt laws for the government of people in futurity. We can, for the time being, point out laws for present necessities. \hl{He supposed there has not yet been a perfect revelation given, because we cannot understand it, yet we receive a little here and a little there. He should not be stumbled if the prophet should translate the bible forty thousand times over and yet it should be different in some places every time, because when God speake, he always speaks according to the capacity of the people.} The starting point for the government of the kingdom is in the Book of Doctrine and Covenants, but he does not know how much more there is in the bosom of the Almighty. When \sout{he} God sees that his people have enlarged upon what he has given us he will give us more. The sta[r]ting point is here, but God has not come here, He has sent his agent, his minister to act in his name. And if he has got an agent to dictate to us here the organization is here. When a man is clothed with authority to do all business for those who sent him, what he does [is?] right, and this is the kind of agent we have got, and God appointed him We did not appoint him...He would rather have the pure revelations of Jesus Christ as they now stand, to carry to the nations, than any thing else. He dont want the legs, nor the head, nor the back bones, nor any thing else of the old horse. We dont <want> to be like the sectarian churches, who, when they have seen errors would turn and run from the truth and run into deeper mud than before.
			
			\item[8 October 1854, CR 100 317] There are a great many persons who are so ancious to learn about eternity, Gods, Angels, heavens, and hells, that they neglect to learn the first lessons preparitory to learning the things they are reaching after. They will come short of them.
		
			I wish to speak a few words about the bible as I have hinted at it. The ordenances of the kingdom of God on <the> earth are the same to the childeren of Adam from the commencement to the end of his posterity pertainig to the carnal state on this earth, and the winding up scene of this mortality. With regard to the Bible we frequently say, We beleive the Bible, but sercumstances alters cases, for what is now required of the <people,> may not be required of a people that may live a hundered years hence. But I wish you to understand, with regard to the ordenances of Gods house to save the people in the Celestial kingdom of <our> God there is no change from the days of Adam to the present time, neither will there be until the last of his posterity is gathered into the kingdom of God. \textit{(I transcribed this quote myself from the longhand report on 8/12/2021; see the manuscript \href{https://catalog.churchofjesuschrist.org/assets/65cd9470-c5b4-4922-8e5c-bf0eeb976187/0/0}{here}.)}
			
			\item[13 July 1862, JD 9:311] When God speaks to the people, he does it in a manner to suit their circumstances and capacities. He spoke to the children of Jacob through Moses, as a blind, stiff-necked people, and when Jesus and his Apostles came they talked with the Jews as a benighted, wicked, selfish people. They would not receive the Gospel, though presented to them by the Son of God in all its righteousness, beauty and glory. Should the Lord Almighty send an angel to re-write the Bible, it would in many places be very different from what it now is. And I will even venture to say that if the Book of Mormon were now to be re-written, in many instances it would materially differ from the present translation. According as people are willing to receive the things of God, so the heavens send forth their blessings. If the people are stiff-necked, the Lord can tell them but little.
			
			\item[7 September 1864] I say to these Latter-day Saints some things I can tell you some things I can't tell you those things that you can't understand and will not use prudently I had better keep them myself the saints have got to be taught according to them they receive and improve upon but tell them a great many things that yet will be told the saints would be like casting pearls before swine this I don't like to do but we want the people who have received the gospel and wish to become saints indeed to receive a little here a little there and improve upon themselves little until they are perfection that is what we want and this is our labor as to the circular and what is in it it is right and still there is things in it that people don't understand [space] there is some things were not written they should understand [space] and I say to you as I can to the whole world there has not yet been a revelation given to this people from the time Joseph commenced to receive revelation but what would be [altered/ltrd?] provided the people were capable of receiving more [space] the Lord has to speak to the people according to their capacity not according to his capacity we are not prepared to receive all the heavens has for us the Lord gives revelation upon revelation here a little and there a little [space] those precepts he gives we should improve upon them as for selling of grain we take measures inasmuch as we can to have the people keep the grain still a great deal will be sold and the men that doesn't feel to sustain what the convention dictated they will sooner or later manifest their folly by going to the dark more and more more and more [sic] until they [love/lv?] the world and every project that is proposed by the leaders of this people and every plan adopted for the benefit of the people those that doesn't hear it will dwindle away sooner or later and finally will apostatize [space] I tell you some things in the circular meant to be in there that I understand you don't [space] I shall not tell it \textit{(This is from a transcript LaJean made of Watt shorthand, a record of BY's remarks on a Southern tour; see online \href{https://catalog.churchofjesuschrist.org/assets/92e3ef48-c3c0-4d38-85d9-96e90d2b4fc6/0/0}{here}.)}
			
		\end{compactdesc}
		
	% -----
	\subsubsection{Harold B. Lee} % *LEE *HAROLDBLEE
	% -----
		\label{EPIST-REVELATION-LEE}
	
		\begin{compactdesc}
			\item[{\href{https://www.churchofjesuschrist.org/study/ensign/1971/01/gods-kingdom-a-kingdom-of-order?lang=eng}{January 1971}}] Keep in mind that the principles of the gospel of Jesus Christ are divine. Nobody changes the principles and doctrines of the Church except the Lord by revelation. But methods change as the inspired direction comes to those who preside at a given time.
		\end{compactdesc}
		
	% -----
	\subsubsection{Russell M. Nelson} % *NELSON *RUSSELLMNELSON
	% -----
		\label{EPIST-REVELATION-NELSON}
	
		\begin{compactdesc}
			\item[{\href{https://www.churchofjesuschrist.org/study/general-conference/2021/10/47nelson?lang=eng}{October 2021}}] Under the Lord's direction and in answer to our prayers, recent procedural adjustments have been made. \textit{He} is the One who wants you to understand with great clarity exactly what you are making covenants to do. \textit{He} is the One who wants you to experience fully \textit{His} sacred ordinances. \textit{He} wants you to comprehend your privileges, promises, and responsibilities. \textit{He} wants you to have spiritual insights and awakenings you've never had before. This He desires for \textit{all} temple patrons, no matter where they live.
			
			Current adjustments in temple procedures, and others that will follow, are continuing evidence that the Lord is actively directing His Church. He is providing opportunities for each of us to bolster our spiritual foundations more effectively by centering our lives on Him and on the ordinances and covenants of His temple. When you bring your temple recommend, a contrite heart, and a seeking mind to the Lord's house of learning, \textit{He} will teach you.
		\end{compactdesc}
		
	% -----
	\subsubsection{Augustine}
	% -----
		\label{EPIST-REVELATION-AUGUSTINE}
		
			\begin{compactdesc}
				\item[DDC I:22--23, 27] Let us consider this process of cleansing as a trek, or a voyage, to our homeland; though progress towards the one who is ever present is not made through space, but through goodness of purpose and character. 23. This we would be unable to do, if wisdom itself had not deigned to adapt itself to our great weakness and offered us a pattern for living; and it has done so actually in human form because we too are human. But because we act wisely when we come to wisdom, wisdom has been thought by arrogant people to have somehow acted foolishly when it came to us; and because we recover strength when we come to wisdom, wisdom has been reckoned as being somehow weak when it came to us. But `the foolishness of God is wiser than men, and the weakness of God is stronger than men'. So although it is actually our homeland, it has also made itself the road to our homeland... The way to health is through medical care; God's care has taken it upon itself to heal and restore sinners by the same methods. When doctors bind wounds, they do this not just anyhow, but in an appropriate manner, so that the effectiveness of the ligature is matched by a kind of beauty; similarly the treatment given by wisdom was adapted to our wounds by its acceptance of human nature, healing sometimes by the principle of contrariety, sometimes by that of similarity.
				\item[DDC I:88] Anyone with an interpretation of the scriptures that differs from that of the writer is misled, but not because the scriptures are lying. If, as I began by saying, he is misled by an idea of the kind that builds up love, which is the end of the commandment, he is misled in the same way as a walker who leaves his path by mistake but reaches the destination to which the path leads by going through a field. But he must be put right and shown how it is more useful not to leave the path, in case the habit of deviating should force him to go astray or even adrift. 
			\end{compactdesc}
	
% ----------------------
\subsection{Commentary and Studies} % *COMMENTARY *STUDIES
% ----------------------
	\label{EPIST-REVELATION-COMMENTARY}

	% -----
	\subsubsection{Principles of revelation}
	% -----
		\label{EPIST-REVELATION-principles}
	
		\begin{compactitem}
			\item God speaks to us ``according to our language'' (\hyperref[2NEPHI-31-3]{2 Nephi 31:3}), where a ``language'' comprises a conceptual scheme and an actual mode of communication (\hyperref[1NEPHI-1-2]{1 Nephi 1:2}).
				\begin{compactitem}
					\item From chapter 1 of \textit{All Things Shining}, by Hubert Dreyfuss and Sean Kelly---helpful for thinking about how far the ``language'' goes into culture and concepts: ``The way of life of a culture is not an explicit set of beliefs held by the people living in it. It is much deeper than that. A person brought up in a culture learns its way of life the way he learns to speak in the language and with the accent of his family and peers. But a way of life is much broader than this. It involves a sense for how it is appropriate and inappropriate to act in each of the social situations one normally encounters; a familiarity with how to make sense of things and of how to act in the everyday world; and most general of all, a style, such as aggressive or nurturing, that governs the actions of the people in the culture although they are normally not aware of it. We can think of it as a cultural commitment that, to govern people’s behavior, must remain in the background, unnoticed but pervasive and real.''
				\end{compactitem}
		\end{compactitem}